% v3.7 (2026-07-17): finite-basis scope closure (conditional joint-extension general theorem; exact atom-cover Engine witness); operational-to-channel, smooth-response-factor, and physical-substrate-anchor gates typed; response-orbit quotient gauge separated from the physical anchor.
%   Aligns the main paper with its self-contained Technical Supplement; internalizes the finite
%   operational basis and finite minimal-realization results as the supplement upstream foundation;
%   adds a theorem-provenance ledger and an explicit division between APF-native, imported-standard,
%   and hybrid results. No downstream geometry claim is strengthened in this version.
%
% v3.5 (2026-07-17): exact closure audit of the remaining record-channel obligations.
%   Derives countability of the physical channel family from finite protocol specification;
%   replaces the impossible tower-faithful continuum code by quotient-faithfulness after
%   asymptotically null refinement histories are identified; proves that a faithful code of the
%   unreduced inverse-limit tower remains zero-dimensional and therefore cannot itself generate
%   a connected continuum; and isolates the two genuinely remaining APF theorems: finite
%   closed-interface generation and finite smooth generation.  No extra basis or smoothness
%   assumption is hidden in the topology.
%
% v3.5 (2026-07-17): topology correction, channel-basis hardening, and theorem-scope audit.
%   Replaces finite-valued terminal records by finite-stage precision-refinement channels;
%   rejects the clopen finite-stage topology, constructs the record topology from faithful real
%   channel limits, and derives finite-channel bounds from additive standing-capacity
%   charges rather than from finiteness of outcome sets; rewrites finite differential
%   presentation in terms of finitely many real channel limits; and states the imported orbit
%   theorem with the source-level completeness condition (closure under pushforwards by local
%   flows).  This removes the finite-image obstruction identified by the v3.2 hostile audit.
%
% v3.2 (2026-07-17): differential-space strengthening and consolidation.
%   Separates executable record observables from their mathematically saturated smooth functional
%   calculus; replaces the overstated finite-generation argument by a precise finite differential
%   presentation criterion; corrects the derivation--cotangent pairing to a canonical isomorphism;
%   states the imported subcartesian orbit theorem with its local-completeness and immersion scope;
%   and consolidates the topology--sheaf--orbit--kernel chain into one typed theorem spine.
%
% v3.1 (2026-07-17): comparison-algebra and differential-space reconstruction.
%   Replaces the circular first-jet/manifold-first construction by a record topology and a
%   refinement-complete sheaf of physical readouts. Derives Sikorski smooth-functional closure
%   from lossless scalar postprocessing, derives local finite generation from finite capacity,
%   and imports the subcartesian orbit theorem so smooth manifolds arise as continuation orbits.
%   Derivation tangents are primary; m/m^2 is retained as an algebraic cotangent and singularity
%   detector. The kernel-first anchor is then recovered on smooth constant-rank continuation orbits.
%
% v3.0 (2026-07-17): primitive record-basis existence and closure of the smooth-certificate gate.
%   Derives a finite, distinction-complete, irredundant first-order record basis from the APF
%   positive cost floor, finite local capacity, closed-world completeness, and lossless minimality.
%   Defines the operational first-order response space as the dual of regularized record germs
%   modulo first-order nullity, proving dimension-completeness and nondegeneracy by construction.
%   Local refinement stability then integrates the basis into a record-coordinate certificate,
%   closing existence of the maximal smooth branch throughout geometric Regime R.
%
% v2.9 (2026-07-17): record-atlas existence, maximal regular locus, and derived first-jet regularity.
%   Replaces the informal smooth-branch clause by a record-atlas certificate: finite families of
%   physically realizable record readouts with full-rank differential generate local charts;
%   compatible overlaps generate a unique maximal C^r atlas.  The regular response locus is
%   proved open and maximal, its rank strata support smooth contextual quotient bundles, and
%   first-jet regularity follows from the record coframe rather than from global injectivity.
%
% v2.8 (2026-07-17): derivation of first-jet regularity on the smooth geometric branch.
%   Defines the intended smooth branch operationally by record-generated local charts on the
%   reduced physical response carrier. Closed-world completeness forces every chart-coordinate
%   distinction into the admitted comparison family; the coordinate differentials form a
%   cotangent coframe, hence the total readout differential is jointly separating. First-jet
%   regularity and the tangent-kernel theorem are therefore derived rather than assumed.
%
% v2.7 (2026-07-17): closed-world germ completeness and regular-stratum kernel theorem.
%   Replaces the previously circular first-order completeness principle by a two-level
%   derivation: closed-world ledger closure first proves equality of physical-response and
%   comparison fibers at the level of continuation germs; first-jet regularity then descends
%   that exactness to tangent kernels. Singular points are treated by finite jets rather than
%   silently assuming derivative faithfulness. Pointed linear readouts and the reduced physical
%   response carrier are constructed explicitly.
%
% v2.6 (2026-07-17): quotient-typed anchor and held-profile repair.
%   Repairs the comparison-response anchor by inserting the held-to-response quotient
%   differential, defines positivity on contextual tangent classes, and distinguishes the
%   invariant record quotient from outcome-bearing collapse instruments.
%
% v2.1 (2026-07-16): held-profile and corpus type alignment.
%   Imports the mature Paper 10 held-profile language; separates held carriers, process
%   algebras, response quotients, records, positive response geometry, and signed causal
%   geometry; distinguishes quantum, record, and causal orientations; strengthens the
%   comparison anchor by contextual identity completeness; and moves the weight-one
%   mechanism to a candidate-realization appendix.
%
% v2.0 (2026-07-16): response-geometry reconstruction.  The paper now separates
%   process geometry, quotient response geometry, and substrate comparison geometry;
%   derives the pullback Finsler gauge and intrinsic response distance; recovers the
%   finite linear quotient theorem of Paper 2 as a flat specialization; separates
%   path-length cost from quadratic residual energy; introduces the comparison-response
%   anchor; prevents a positive pullback metric from being conflated with Lorentzian
%   comparison signature; adds the trace-volume theorem and the operator target
%   P_tr K^{-1}J=0 for the no-trace weak-field gate; and rewrites the closure roadmap.
\documentclass[11pt]{article}

\usepackage[utf8]{inputenc}
\usepackage[T1]{fontenc}
\usepackage{lmodern}
\usepackage{microtype}
\usepackage[letterpaper,margin=1in]{geometry}
\geometry{margin=1in}
\setlength{\emergencystretch}{4em}
\usepackage{amsmath,amssymb,amsthm,mathtools,mathrsfs}
\usepackage{booktabs,tabularx,array,longtable}
\usepackage{enumitem}
\usepackage{xcolor}
\usepackage[colorlinks=true,linkcolor=black,citecolor=black,urlcolor=black]{hyperref}
\usepackage[nameinlink,capitalize]{cleveref}

\hypersetup{
  colorlinks=true,
  linkcolor=black,
  citecolor=black,
  urlcolor=black,
  pdfauthor={E. S. Brooke},
  pdftitle={Response-Induced Geometry of the Admissibility Substrate}
}

\setlist[itemize]{leftmargin=1.4em, topsep=3pt, itemsep=2pt}
\setlist[enumerate]{leftmargin=1.6em, topsep=3pt, itemsep=2pt}
\renewcommand{\arraystretch}{1.15}

\newtheorem{theorem}{Theorem}[section]
\newtheorem{lemma}[theorem]{Lemma}
\newtheorem{proposition}[theorem]{Proposition}
\newtheorem{corollary}[theorem]{Corollary}
\newtheorem{definition}[theorem]{Definition}
\newtheorem{assumption}[theorem]{Assumption}
\newtheorem{principle}[theorem]{Principle}
\newtheorem{gate}[theorem]{Gate}
\theoremstyle{remark}
\newtheorem{remark}[theorem]{Remark}
\newtheorem{warning}[theorem]{Boundary}
\newtheorem{audit}[theorem]{Audit Note}

\newcommand{\G}{\Gamma}
\newcommand{\Ccap}{C}
\newcommand{\entails}{\Vdash}
\newcommand{\goes}{\rightsquigarrow}
\newcommand{\Cost}{\kappa}
\newcommand{\Cont}{\operatorname{Cont}}
\newcommand{\Erec}{E_{\rm rec}}
\newcommand{\phide}{\phi_{D\to E}}
\newcommand{\Req}{\phi_{\rm eq}}
\newcommand{\Kgeom}{\mathcal K}
\newcommand{\Mgeom}{\mathcal M}
\newcommand{\Dgeom}{\mathcal D}
\newcommand{\Egeom}{E_{\rm geom}}
\newcommand{\Jsrc}{\mathcal J}
\newcommand{\Status}[1]{\textsf{#1}}
\newcommand{\R}{\mathbb R}
\newcommand{\eps}{\varepsilon}
\newcommand{\dvol}{\,d^3x}
\newcommand{\Boxop}{\square}
\newcommand{\Lie}{\mathcal L}
\newcommand{\Pproc}{\mathcal P}
\newcommand{\Qresp}{\mathcal Q}
\newcommand{\Fresp}{F_{\rm resp}}
\newcommand{\dresp}{d_{\rm resp}}
\newcommand{\dproc}{d_{\rm proc}}
\newcommand{\Iphys}{I_{\rm phys}}
\newcommand{\Gcost}{g^{\rm cost}}
\newcommand{\Qcmp}{Q^{\rm cmp}}
\newcommand{\Ptr}{P_{\rm tr}}
\newcommand{\Ptf}{P_{\rm TF}}
\newcommand{\Held}{\mathsf{Held}}
\newcommand{\Rec}{\mathsf{Rec}}
\newcommand{\VH}{V_H}
\newcommand{\Tproc}{\mathfrak T_H}
\newcommand{\AH}{\mathcal A_H}
\newcommand{\Qrec}{\mathsf Q}

\title{Response-Induced Geometry of the Admissibility Substrate\\
\large Quotient Response Geometry, Intrinsic Cost, and the Weak-Field Trace Obstruction}
\author{E.\ S.\ Brooke \\[0.4em] \normalsize ORCID: \href{https://orcid.org/0009-0001-2261-4682}{0009-0001-2261-4682}}
\date{Version 3.7 --- July 2026}

\usepackage{seqsplit}%CORPUSGUARD
\newcommand{\codeid}[1]{\texttt{\seqsplit{#1}}}%CORPUSGUARD
\newcommand{\coderef}[2]{{\small(\codeid{#1}, \codeid{#2})}}%CORPUSGUARD

%CORPUSFMT 2026-07-04: preamble normalized to the Paper 0 corpus standard
\usepackage{parskip}
\usepackage[most]{tcolorbox}
\usepackage{titlesec}
\titleformat{\section}{\Large\bfseries}{\thesection.}{0.5em}{}
\titleformat{\subsection}{\large\bfseries}{\thesubsection}{0.5em}{}
\titleformat{\subsubsection}{\normalsize\bfseries}{\thesubsubsection}{0.5em}{}
\setlength{\emergencystretch}{1.5em}

\begin{document}
\maketitle

\begin{abstract}
This paper develops the geometric-substrate program of the Admissibility Physics Framework as a response-geometry program.  Its upstream object is now the structurally aligned held continuation profile of Paper~10: a complete pre-record carrier whose contextual identity is preserved through admissible coherent continuation.  The central type distinction is among the held-profile carrier $\VH$, its represented process algebra $\Tproc$, a finite response-coordinate space and response-null quotient, the invariant committed-record image $\mathsf R(H)$ and the outcome-bearing collapse maps $\{\mathsf C_i\}$, and the physical comparison substrate.  The represented process algebra supplies the operational target of the response map and keeps held continuation geometry, response geometry, record formation, and substrate geometry distinct.  Once redundant response coordinates are quotiented, a differentiable response map pulls the operational process norm back to a positive homogeneous tangent gauge
\[
        \Fresp(q,v)
        =\left\|D\bar\Phi_q(v)\right\|_{\rm proc}.
\]
This gauge defines an intrinsic response length and path distance.  In a finite linear chart it is flat and translation invariant, and the resulting distance theorem is exactly the geometric form of the finite quotient, upper-gain, minimum-gain, and bi-Lipschitz package developed in Paper~2.

The response distance is not automatically the physical interaction ledger.  This paper proves the precise calibration theorem: if the physical ledger rate is bounded above and below by the pullback gauge, then the two induced path costs are bi-Lipschitz equivalent, with equality only under exact rate calibration.  It also separates one-homogeneous path cost from the two-homogeneous residual energy of smooth residual realization.  In the latter branch the local theorem is
\[
        I_{\Pi}^{\rm res}(q)
        =\frac12 d_{g^{\rm E}}(q_0,q)^2+o(d_{g^{\rm E}}^2),
\]
not an identification of interaction energy with distance itself.

The intrinsic response-orbit quotient is constructed kernel-first, but no manifold and no physical substrate anchor are assumed upstream.  Finite-stage executable tests are organized into precision-refinement channels: each stage has a finite outcome set, while the compatible tower may carry a real-valued limit.  A conditional finite operational-basis theorem controls operational distinction content; a separate faithful-realization gate is required before that basis becomes a jointly realizable, state-separating precision-channel family.  Quotient-faithful channel limits generate the record topology, while second countability additionally requires a countable specification grammar and a countable topology-generating subfamily.  Lossless relabelling preserves the physical quotient and smooth saturation supplies the minimal Sikorski functional calculus.  The resulting object is a differential space with derivation tangent
\[
T_p^{\rm der}P=\operatorname{Der}_p(\mathscr C_p,\mathbb R),
\]
canonically dual to $\mathfrak m_p/\mathfrak m_p^2$.  Finite smooth generation is isolated as a separate subcartesian gate; on that locus the established orbit theorem makes complete-family continuation orbits immersed smooth manifolds.  Rank-changing, non-flowable, or non-finitely-presented sectors remain intrinsic singular strata.

On a smooth continuation orbit, set-level closed-world fibre completeness is not enough for first-order faithfulness.  If the complete response and a finite local APF readout generator factor smoothly through one another, then
\[
\ker D\beta_x
= N_x^{\rm ctx}
=\bigcap_{a\in\mathcal A_x}\ker\ell_{a,x},
\]
and the induced quotient response is injective.  Its process norm then defines an intrinsic positive response-orbit gauge.  Transport of that gauge to physical site displacements remains a separate substrate-anchor gate.  Quadraticity may convert the response gauge into a positive cost metric, while Lorentzian comparison geometry still requires an independent causal cone and signed two-homogeneous structure.

The weak-field program is retained and sharpened.  In each clock--ruler comparison plane the static load response has generator $L$.  If the loaded comparison kernel is $q_\chi=R_\chi^*q_0$, then
\[
        \left.\frac{d}{d\chi}
        \log\sqrt{|\det q_\chi|}\right|_{\chi=0}
        =\operatorname{tr}L.
\]
After clock calibration this scalar volume mode is exactly $1-\gamma_C$.  The remaining APF-specific theorem target is therefore the operator statement
\[
        \Ptr\mathcal K^{-1}\mathcal J=0,
\]
where $\mathcal K$ is the gauge-fixed response Hessian and $\mathcal J$ the scalar recruitment source.  It states that static scalar load excites only trace-free exchange and no independent comparison-volume mode.  Paper~6 remains the source of record for the represented smooth GR closure; Paper~9 now supplies the intrinsic response geometry beneath that closure and identifies its weak-field scalar obstruction.
\end{abstract}


\begin{tcolorbox}[title={Companion technical supplement and closure policy},colback=gray!4,colframe=black!55]
This main paper is accompanied by \emph{Paper 9 Technical Supplement:
Dependency-Closed Conditional Package for Response-Induced Geometry of the
Admissibility Substrate}, Version 1.1.  The supplement reproduces the premises
used by the certified native core and keeps finite joint extension, compatible
joint realization, operational-to-channel realization, countable topology
generation, finite smooth generation, smooth response-factor completeness, the
physical substrate anchor, causal signature, no trace, and nonlinear GR as
named gates.  No external APF reading is required for that certified native
core; Paper~6 remains the separately marked external source of record for the
represented nonlinear endpoint.
\end{tcolorbox}

\section*{Result provenance classes}
Every load-bearing result in this paper belongs to one of three classes.
\textbf{APF-native} results are proved from locally printed premises and any
explicitly named APF regime gates.  \textbf{Standard imports} are classical
mathematical theorems stated with their exact hypotheses.  \textbf{Hybrid}
results combine an APF certificate with a standard theorem.  Open gates and
external endpoints are not promoted into the APF-native class by notation.

\begin{center}
\fbox{\begin{minipage}{0.86\textwidth}
\centering
\vspace{3pt}
\textbf{Governing claim.}\\[3pt]
Geometry is the represented intrinsic structure of held-profile response, after contextual identity, null response directions, record formation, ledger calibration, substrate anchoring, quadraticity, and causal signature are kept type-distinct.\\[4pt]
{\small
\[
\begin{aligned}
\mathscr A_{\rm fin}
&\longrightarrow (P^{\rm red},\tau_{\rm rec},\mathscr C_{\rm APF}) \\
&\longrightarrow \text{finite differential presentation} \\
&\longrightarrow \text{continuation orbits}
\longrightarrow T S/N^{\rm ctx} \\
&\longrightarrow \Fresp
\longrightarrow \dresp
\longrightarrow \text{calibrated cost geometry} \\
&\longrightarrow \text{signed causal geometry}
\longrightarrow \text{volume response.}
\end{aligned}
\]
}
\vspace{3pt}
\end{minipage}}
\end{center}

\tableofcontents

\section{Orientation: the geometric-substrate problem}

Paper 10 supplied a continuation-first language for APF.  Its central move was to separate admissible possibility from dynamics:
\[
        \G\entails_C D\goes E
\]
states that \(D\) can continue as \(E\) within capacity bound \(C\).  A dynamics, probability law, or variational rule is a later selector or weighting on admissible continuation words.

The same paper proposed a compressed reading of geometry:
\begin{quote}
Geometry is the regularized cost structure of carrying comparison, localization, and transport continuations.
\end{quote}
This paper begins the work required to make that sentence theorem-grade.

The geometric-substrate problem has six type-distinct parts:
\begin{enumerate}
\item define comparison continuations and finite response models operationally;
\item quotient response-null coordinates and construct intrinsic response length;
\item calibrate that modeled length to the physical interaction ledger;
\item anchor response directions to tangent directions of the comparison substrate;
\item derive a signed causal comparison geometry without confusing it with the positive cost metric;
\item determine how recruitment load deforms the comparison geometry and align the result with the Paper~6 represented-GR closure.
\end{enumerate}

This revision proves the finite geometric statements in items 2 and the conditional calibration statement in item 3.  It formulates item 4 as an explicit anchor gate, corrects the signature logic in item 5, and retains the redshift and scalar trace program as the first worked application of item 6.

\begin{warning}[Status discipline]
Throughout this paper, \Status{P} means proved from stated definitions and assumptions in this paper; \Status{C} means conditional on an explicit gate or structural commitment; \Status{R} means recovered in a represented regime; \Status{I} means imported from another APF paper under its stated hypotheses; and \Status{O} means open.  The present document mixes \Status{P}, \Status{C}, \Status{R}, and \Status{I}: the quotient and pullback theorems are finite mathematics; ledger realization, anchoring, signature, and the no-trace source are conditional or open; redshift is a conditional recovery; and the Einstein closure is imported from Paper~6.
\end{warning}

\subsection*{On operational language}

This paper develops geometry as the represented cost structure of finite comparison continuations.  The continuation primitive imported from Paper~10 is a static algebraic relation $\G\entails_C D\goes E$ on continuation profiles, and the metric-emergence ladder of \S\ref{sec:metric-emergence} derives the quadratic metric kernel as a representation \emph{conditional on a series of structural gates closing simultaneously} --- not as something the framework's primitives produce by temporal evolution.

The phrasings ``becomes available,'' ``emerges,'' ``closes,'' and ``recovers'' that appear throughout this paper are the local readings of conditional structural commitments.  Where the prose says ``the metric becomes available only after several gates close,'' the structural reading is ``the metric is a representation conditional on the gates' joint structural validity.''  Where the prose says ``admissible chart changes become diffeomorphism-type relabelings in the smooth represented regime,'' the structural reading is ``in the smooth represented regime, the admissible relabelings are exactly the diffeomorphism-type ones.''  The ``becomes'' is locator-language, not process-language.

Time itself is derived in this corpus: Paper~3 derives the arrow of time as the loss of global admissible coordination, and Paper~6 derives the metric and dimensionality of spacetime.  The geometric machinery developed here operates on the static comparison structure $(c_x, q_x)$ at the substrate; what reads as ``response,'' ``recruitment,'' or ``redshift recovery'' is the local reading of conditional structural commitments under temporal slicing.  A reader who wants the structural reading can take every ``becomes available,'' ``emerges,'' and ``recovers'' as shorthand for the local reading of the corresponding structural conditional.  The technical apparatus developed below --- the cost germ $c_x$, the quadratic kernel $q_x$, the geometric recruitment functional, the cost-curvature response --- is already structural; only the pedagogical wrapping is temporal.

This convention matches Paper~0's ``A note on language'' in the Prelude \S1 and Paper~0's Descriptive Reading chapter on temporal language and the eternalist commitment, Paper~1 supplement v8.24's ``On the use of operational language'' note in \S1, and Paper~10 v1.17's parallel note in \S1.  The framework is in this respect in the company of block-universe readings of general relativity, the Wheeler--DeWitt timeless quantum gravity programme, and the Page--Wootters mechanism for time emergence in entangled quantum systems --- structural at the foundation, with time as a derived feature of how the structure is read locally.

\smallskip\noindent\textit{A note on two ``substrate'' senses.}\quad The framework uses ``substrate'' in two distinct but related senses, and a reader of this paper alongside Paper~0 should keep them apart. Paper~9 uses \emph{substrate} in its \textbf{material-substrate} sense throughout: the \emph{geometric substrate} is the physical object that carries comparison continuations; the \emph{comparison substrate} is the physical object equipped with admissible comparison continuations between sites; \emph{substrate degrees of freedom} are the local mode-content that the recruitment profile $\varphi_d(x)$ recruits to carry a continuation. Paper~0 \S\texttt{sec:substrate\_position\_axis} introduces a different sense: the \emph{substrate-position axis} is a position locus on the framework's three-regime architecture (void regime / middle regime / saturation regime). The two senses are parallel but not identical: the geometric substrate is \emph{what carries} the comparison continuation; the substrate-position axis is \emph{where in regime space} the architecture sits. Paper~9's content is largely middle-regime physics in the substrate-position-axis sense (the metric-emergence ladder of \S\ref{sec:metric-emergence} requires middle-regime conditions: locally smooth and locally additive comparison cost on a connected comparison substrate without saturation); the saturation-direction reading is articulated in \S``Cost-curvature response'' where recruitment-load already committed reduces residual stiffness for additional comparison continuations. Paper~6 main v3.9 + Supplement v2.10 carry the substrate-position-axis architectural reading at the dynamics/geometry/cosmology level.


\section{Canonical type ledger and corpus dependencies}\label{sec:type-ledger}

Paper~9 consumes the mature continuation architecture rather than redefining it.  The canonical objects are:
\begin{longtable}{p{0.17\textwidth}p{0.49\textwidth}p{0.23\textwidth}}
\toprule
Object & Meaning & Source of record \\
\midrule
$H\in\Held$ & Structurally aligned, occupied, pre-record continuation profile & Papers~1, 10, and 14 \\
$\VH$ & Complete finite represented carrier of held continuation profiles & Papers~1 and 10; Paper~5 supplies the quantum-sector completion \\
$\Tproc$ & Represented algebra of admissible transformations of held profiles & Papers~1 and 5 \\
$\AH$ & Observable/effect algebra paired with held states and processes & Paper~5 \\
$\mathsf R(H)\in\mathsf{InvRec}$ & Invariant phase-insensitive committed-record image & Papers~10 and 37 \\
$\{\mathsf C_i(H)\}$ & Outcome-bearing committed records & Papers~5 and 37 \\
$R_\Pi/\!\sim$ & Response coordinates modulo process-null implementation directions & Paper~9 \\
$\Mgeom$ & Physical comparison substrate carrying localization and transport continuations & Paper~9 \\
\bottomrule
\end{longtable}

The geometry pipeline of this paper is therefore
\[
\begin{aligned}
H\in\Held&\longrightarrow \VH\longrightarrow\Tproc,\\
T\Mgeom&\longrightarrow R_\Pi/\!\sim\longrightarrow\Tproc.
\end{aligned}
\]
with the final arrow supplied only by a physical comparison-response anchor.  Record formation $\Qrec:\Held\to\Rec$ is downstream and is not used to define the pre-record response geometry.

\begin{audit}[Three distinct planes and three distinct orientations]
The elementary quantum comparison carrier of Paper~5, the clock--ruler load-response plane of this paper, and the tangent plane of a Lorentzian comparison substrate are not automatically identical:
\[
W_A^{\rm quantum}\neq W_e^{\rm clock\text{-}ruler}\neq T_x\Mgeom
\]
unless an explicit anchor theorem identifies them.  Likewise distinguish
\[
\begin{aligned}
o_H&=\text{held complex orientation},\\
o_{\rm rec}&=\text{irreversible record order},\\
o_{\rm caus}&=\text{time orientation of the causal cone}.
\end{aligned}
\]
Paper~5 derives $o_H$ in its declared coherent quantum regime.  Papers~3 and 37 supply the irreversible record order $o_{\rm rec}$.  Paper~9 may represent that order by $o_{\rm caus}$ only after the causal-cone and signed-signature gates close.  No one of these orientations is substituted for another by notation.
\end{audit}

\begin{audit}[Cross-paper dependency boundary]
Paper~10 supplies continuation grammar, held-profile identity, and the universal held continuation lift.  Paper~5 supplies the quantum-sector represented carrier, local coherent circle, and natural complex orientation.  Paper~9 neither rederives those quantum results nor uses them to close the substrate anchor, the Lorentzian signature gate, or the weak-field no-trace operator.  Paper~6 remains the source of record for the represented nonlinear Einstein closure.
\end{audit}

\section{Main results and non-claims of this draft}
\label{sec:main-results}

The paper now has two linked theorem spines.  The first is intrinsic response geometry; the second is the weak-field deformation of the represented comparison kernel.  The table records what is proved, what is conditional, and what remains open.

\begin{longtable}{p{0.25\textwidth}p{0.14\textwidth}p{0.50\textwidth}}
\toprule
Result & Status & Meaning \\
\midrule
Finite operational basis & \Status{Cond./model} & General basis existence is conditional on the joint-realizability package; APF-Engine certifies the exact finite atom-cover instantiation and one compatible minimum joint carrier. \\
Finite physical precision-channel generation & \Status{G} & Requires a faithful realization preserving joint realizability and recoverability, separating reduced states, and admitting quotient-faithful real limits. \\
Record topology & \Status{P$\mid$G} & Real channel limits generate the initial topology and a separating family gives Hausdorffness.  Second countability additionally requires a countable grammar and countable topology-generating subfamily. \\
Differential-space closure & \Status{P/D} & Localization and smooth saturation define the minimal Sikorski structure compatible with the admitted record algebra. \\
Finite-channel sufficiency & \Status{P/G} & Once a complete jointly realizable channel basis exists, additive standing-capacity accounting makes it finite.  Existence follows only after the operational-to-channel gate closes. \\
Finite differential sufficiency & \Status{G} & Finite state separation does not imply finite smooth generation of the continuation-regular sheaf. \\
Smooth continuation orbits & \Status{I/P} & On the finite-differential-presentation locus, the imported subcartesian orbit theorem gives immersed smooth continuation orbits for complete local flow families. \\
Regular local response quotient & \Status{P$\mid$H} & Constant rank gives quotient charts; derivative-kernel faithfulness additionally requires smooth two-way response factorization. \\
Response-orbit quotient gauge & \Status{P$\mid$H} & The injective quotient response defines an intrinsic positive homogeneous gauge on the response orbit after the smooth-factor gate closes. \\
Physical comparison-substrate anchor & \Status{G} & Requires a smooth site map and tangent map equal to its quotient differential, with direct-comparison calibration and stable provenance. \\
Intrinsic response distance & \Status{P} & Integrating the response-orbit gauge gives a coordinate-invariant path distance and endpoint process bound. \\
Linear Paper~2 specialization & \Status{P} & The finite linear branch is the flat quotient $R_\Pi/\ker T_\Pi$ and recovers the optimal response gains. \\
Physical-ledger calibration & \Status{P$\mid$H} & Pointwise comparison of ledger rate with the response gauge gives bi-Lipschitz comparison; equality requires exact calibration. \\
Local residual energy--distance law & \Status{P$\mid$SRR} & A smooth positive Hessian gives $I_\Pi^{\rm res}=\tfrac12 d_{g^{\rm E}}^2+o(d^2)$ near equilibrium. \\
Quadratic positive cost metric & \Status{P$\mid$Q} & A parallelogram certificate converts the response gauge into a positive metric; otherwise the Finsler branch remains. \\
Lorentzian comparison structure & \Status{O/C} & Requires an independent finite-coordination cone and signed two-homogeneous structure. \\
Point-source exterior profile & \Status{C/P} & An assumed isotropic quadratic exterior kernel gives $\chi(r)\sim1/r$; kernel and calibration remain inputs. \\
Clock-sector redshift recovery & \Status{R} & Calibrated load, clock response, and the proper-time dictionary recover the redshift factor conditionally. \\
Trace--volume theorem & \Status{P/C} & The first comparison-volume variation is $\operatorname{tr}L$, equal after clock calibration to $1-\gamma_C$. \\
No-trace operator reduction & \Status{P/C} & The local weak-field target is $\Ptr\mathcal K^{-1}\mathcal J=0$; deriving it from recruitment remains open. \\
Conditional weak-field exterior & \Status{C/R} & If no trace closes, the first-post-Newtonian Schwarzschild exterior follows conditionally. \\
Full represented GR & \Status{I/C} & External Paper~6 endpoint under H1--H3/A9/Lovelock; finite-continuability translation remains open. \\
\bottomrule
\end{longtable}

\begin{quote}
Finite capacity and the positive marginal floor do not by themselves turn an arbitrary closed operational ledger into a jointly realizable finite basis. The general basis theorem uses an explicit finite joint-extension package. The Engine closes the finite atom-cover instantiation and separately exposes the general bridge as open. A further faithful realization gate is required before the operational basis becomes a finite physical precision-channel basis.
\end{quote}

\begin{audit}[Non-claims]
This paper does not identify every physical interaction cost with a distance, does not derive Lorentzian signature from $T^*T$, does not infer spacetime dimension from a register-count theorem, and does not turn physical divisibility into the mathematical subdivision of a curve.  It does not rederive the Einstein equations, the value of $G$, the dimension $(1+3)$, or the full equivalence principle.  Its new mathematical contribution is the typed record-topology--differential-space--orbit construction, the kernel-first response geometry on emergent smooth branches, and the trace--volume reduction of the weak-field spatial obstruction.
\end{audit}

\section{Alignment with Paper 6: represented-GR source of record}
\label{sec:paper6-alignment}

Paper~6 and its Technical Supplement already contain the represented-geometry route to GR.  In that paper the sequence is:
\[
\begin{aligned}
\text{cost distance}
&\Rightarrow g_{\mu\nu}
\Rightarrow (-,+,+,+),\ d=4 \\
&\Rightarrow \text{capacity-gradient curvature}
\Rightarrow \text{A9 closure} \\
&\Rightarrow \text{Lovelock uniqueness}
\Rightarrow
G_{\mu\nu}+\Lambda g_{\mu\nu}=\kappa T_{\mu\nu}.
\end{aligned}
\]
Paper~9 should therefore not duplicate that proof.  Its role is to express the same represented-geometry closure in the newer finite-continuability language and to identify what that closure looks like in the weak-field scalar sector.  Of the imported hypotheses, H1 (continuum emergence) and H2 (locality of the $\Delta_{\mathrm{geo}}$ closure) are discharged by the Recruitment-Radius papers (Papers~24 and~25) via continuous anchor profiles and $L_{\mathrm{loc}}$ on substrate kernels; the represented-geometry import that remains is H3 together with the A9/Lovelock uniqueness step.

\begin{theorem}[Paper 6 represented-GR import]
\label{thm:paper6-gr-import}
Under the Paper~6 hypotheses H1--H3, together with the A9 closure and Lovelock uniqueness in the represented smooth Lorentzian regime, the admissibility-preserving geometric dynamics are governed by
\[
        G_{\mu\nu}+\Lambda g_{\mu\nu}=\kappa T_{\mu\nu}.
\]
This theorem is imported here as a represented-geometry result, not reproved.
\end{theorem}

\begin{corollary}[Weak-field shadow of the Paper 6 closure]
\label{cor:paper6-weak-field-shadow}
If the Paper~6 represented-GR closure is accepted, then in the static weak-field exterior of a point mass the metric representation has the first-post-Newtonian Schwarzschild form
\[
        ds^2
        =-
        \left(1-\frac{2GM}{rc^2}\right)c^2dt^2
        +
        \left(1+\frac{2GM}{rc^2}\right)d\mathbf x^2
        +O(c^{-4}),
\]
and hence the scalar weak-field response of Paper~9 has
\[
        \gamma_C=1,
        \qquad
        \theta=1-\gamma_C=0.
\]
\end{corollary}

\begin{audit}[What Paper 9 adds beyond importing Paper 6]
The import theorem says that the represented smooth geometry satisfies the Einstein equations once the Paper~6 gates close.  Paper~9 adds a lower-level translation: redshift fixes the clock response, while the entire leading spatial ambiguity is the trace of a local clock-ruler scalar response generator.  Thus the no-trace/unimodular gate is the finite-continuability weak-field face of the Paper~6 A9/Lovelock closure.  If Paper~6 is used, \(\theta=0\) follows in the weak-field represented regime.  If Paper~9 is developed bottom-up, \(\theta=0\) is the next primitive substrate theorem to prove.
\end{audit}

\section{Reader path and theorem spine}
\label{sec:reader-path}

The paper should be read in two passes.  The first pass constructs intrinsic response geometry without presupposing spacetime.  The second pass anchors and deforms that geometry in the weak-field comparison sector.

\begin{center}
\fbox{\begin{minipage}{0.92\textwidth}
\textbf{Publication spine.}
{\small
\[
\begin{gathered}
\text{finite operational comparison content}
\xrightarrow{\text{joint-extension package}}
\text{finite operational basis}\\
\xrightarrow{\text{faithful channel realization}}
\text{finite separating precision channels}\\
\xrightarrow{\text{real limits + countable topology generation}}
\text{Hausdorff second-countable record topology}\\
\xrightarrow{\text{finite smooth generation}}
\text{subcartesian differential space}\\
\xrightarrow{\text{complete local flows}}
\text{smooth continuation orbits}\\
\xrightarrow{\text{smooth response-factor completeness}}
T S/N^{\rm ctx}\hookrightarrow T\mathcal O\\
\Longrightarrow \text{intrinsic response-orbit gauge}\\
\xrightarrow{\text{physical substrate anchor}}
\text{physical comparison geometry}\\
\xrightarrow{\text{quadraticity / signed causal gates}}
\delta\log\sqrt{|\det q|}=\operatorname{tr}L=1-\gamma_C.
\end{gathered}
\]
}
\end{minipage}}
\end{center}

The decisive closures are therefore:
\begin{enumerate}
\item close the faithful operational-to-channel, countable-generation, and
finite-smooth-generation gates;
\item certify complete local continuation flows and smooth response-factor
completeness on the regular branch;
\item calibrate the intrinsic response-orbit gauge to the physical interaction
ledger;
\item construct the physical comparison-substrate anchor and independent causal
cone;
\item derive the loaded-equilibrium Hessian $\mathcal K$, source $\mathcal J$,
and the no-trace condition $\Ptr\mathcal K^{-1}\mathcal J=0$.
\end{enumerate}
The paper now supplies precise conditional theorems and exposes each remaining
bridge separately.  Connections and curvature are deferred until regularity,
smoothness, strong convexity or quadraticity, the physical anchor, and causal
signature have all been certified.

\begin{audit}[Why Paper 9 is the geometric home]
Paper~2 supplies a finite response model and its quotient estimates.  Paper~9 generalizes those estimates into intrinsic path geometry, anchors them to the comparison substrate, separates positive cost geometry from signed causal geometry, and applies the result to recruitment-deformed clock--ruler response.  Paper~2 should therefore import finished Paper~9 theorems later rather than absorb the geometric program itself.
\end{audit}

\section{Imports from finite continuability}

\begin{definition}[Continuation judgment]
The primitive APF judgment is
\[
        \G\entails_C D\goes E.
\]
It means that, in context \(\G\), distinction-expression \(D\) can continue as distinction-expression \(E\) using enforcement capacity no greater than \(C\).
\end{definition}

\begin{definition}[Continuation cost]
The cost of a continuation is
\[
        \Cost_\G(D\goes E)
        :=
        \inf\{C:\G\entails_C D\goes E\}.
\]
An identity-class continuation is a continuation with \(E\in[D]_\G\), where \([D]_\G\) is the tested continuation-equivalence class of \(D\).
\end{definition}


\begin{definition}[Held continuation profile import]
A held continuation profile $H\in\Held$ is the structurally aligned, occupied, pre-record object that carries complete contextual continuation identity together with every consequential coherent distinction still available to later continuation.  Its finite represented carrier is $\VH$.  A record-forming class transition is a generally nonfaithful map
\[
\Qrec:\Held\longrightarrow\Rec
\]
that commits a tested distinction and may remove globally common orientation while retaining consequential relative information.
\end{definition}

\begin{definition}[Universal held continuation lift import]
For a local comparison coordinate $q$ inside $H$, Paper~10 defines the minimum complete local carrier
\[
\mathsf L_q(H)=\mathsf G_q(H)/\!\sim_q,
\]
the quotient of admissible local continuation germs by complete short-context equivalence.  In a smooth operationally first-order branch it is represented by the admitted first-jet locus.  Paper~9 uses this only as an identity-complete local carrier; it does not import the Paper~5 quantum $S^1$ theorem into substrate geometry.
\end{definition}

\begin{definition}[Recruitment profile]
For a continuation \(D\goes E\), the recruitment profile \(\phide(x)\) records where and how strongly substrate degrees of freedom are recruited to carry that continuation.  The associated recruitment cost is denoted
\[
        \Erec[\phi;\G].
\]
Fields are represented equilibrium recruitment profiles.  Radiation is cost discharge from failed adiabatic recruitment.  Geometry is the special case in which the recruited substrate carries comparison, localization, and transport continuations.
\end{definition}

\noindent\textit{Source of the recruitment formalism.}\quad The recruitment profile $\phide$, the recruitment cost $\Erec$, and their dynamics are formalized in the Recruitment-Radius papers (Papers~24 and~25): the recruitment cost functional, the master equation for substrate recruitment, continuous anchor profiles, and the substrate-kernel structure.  The present paper specializes that formalism to the geometric case --- comparison, localization, and transport continuations.  Where the derivations below call for a nonlinear recruitment functional, a substrate kernel, or its Green's function, the source of record is Papers~24/25.

\section{Comparison continuations and the three geometric levels}
\label{sec:three-geometries}

\begin{definition}[Comparison continuation]
A comparison continuation from $x$ to $y$ in context $\G$ is a substrate operation $K:x\goes y$ that carries a probe's held continuation profile from $x$ to $y$, preserves its complete contextual identity class along the way, and establishes endpoint correspondence.  Its ledger may be decomposed as
\[
\Cost_\G(K)
=
\Cost_{\rm transport}(K)
+
\Cost_{\rm identity}(K)
+
\Cost_{\rm corr}(K),
\]
with the understanding that the three terms may share physical support.
\end{definition}

\begin{definition}[Directed comparison cost and comparison substrate]
The directed comparison cost is
\[
        d^+_\G(x,y)
        :=
        \inf\{\Cost_\G(K):K\text{ is an admissible comparison continuation }x\goes y\}.
\]
A comparison substrate $\Mgeom$ is a regular domain of distinguishable comparison sites equipped with admissible nearby comparison continuations.  It is locally comparison-complete at scale $\ell$ when every represented pair below that scale is connected by at least one admissible comparison continuation.
\end{definition}

\begin{definition}[Direct local comparison germ]
Whenever the refinement limit exists, define
\[
        c_x^{\rm dir}(u)
        :=
        \lim_{\epsilon\downarrow0}
        \frac{d^+_\G(x,x+\epsilon u)}{\epsilon}.
\]
This is the direct substrate-side cost germ.  It may be asymmetric, nonsmooth, or degenerate before further gates close.
\end{definition}

The response construction below introduces three distinct geometries.

\begin{longtable}{p{0.24\textwidth}p{0.29\textwidth}p{0.37\textwidth}}
\toprule
Level & Carrier & Role \\
\midrule
Process geometry & Represented held-process algebra $\Tproc$ acting on $\VH$, with operational distance $\dproc$ & Measures distinguishability or displacement between represented processes.  In the Paper~2 channel model, $\|\Delta\|_{\rm proc}=\tfrac12\|\Delta\|_\diamond$. \\
Response geometry & Response-null quotient $\Qresp_{\Pi,\mathcal F}$ & Measures intrinsic length of changing a represented response class while redundant coordinates have been removed. \\
Substrate comparison geometry & Comparison substrate $\Mgeom$ & Measures comparison, localization, and transport directions at physical sites after an anchor identifies them with response directions. \\
\bottomrule
\end{longtable}

\begin{audit}[Type ledger]
The following objects must not be silently identified:
\[
\begin{array}{c|c|c}
\text{object} & \text{domain} & \text{homogeneity/role}\\ \hline
\Fresp(q,v) & T\Qresp & \text{positive one-homogeneous gauge}\\
\dresp(q_0,q_1) & \Qresp\times\Qresp & \text{intrinsic path distance}\\
I_\Pi^{\rm path} & \text{physical protocol paths} & \text{calibrated path ledger}\\
I_\Pi^{\rm res} & \text{endpoint response state} & \text{local two-homogeneous residual energy}\\
M_{\rm sep} & \Pproc & \text{extrinsic distance to a free-process set}\\
\Qcmp_x(u) & T_x\Mgeom & \text{signed causal comparison function}.
\end{array}
\]
A valid theorem may compare these objects, but their equality is never a consequence of notation.
\end{audit}

\section{Response-null quotient and intrinsic process geometry}
\label{sec:metric-emergence}

Let $R_\Pi$ be a finite-dimensional response-coordinate manifold, fix a free represented held process $\mathcal F\in\Tproc$, and let
\[
        \Phi_{\mathcal F}:R_\Pi\longrightarrow\Tproc,
        \qquad
        \Phi_{\mathcal F}(0)=\mathcal F,
\]
be a differentiable response model.  Coordinates can represent memory occupation, a residual correlation register, a joint repair operation, nonfactorizing control amplitude, or any other finite response variable.  Coordinates are not physical merely because they are written down.

\begin{definition}[Response fiber and regular quotient chart]
For fixed $\mathcal F$, define
\[
        r\sim_{\mathcal F}r'
        \quad\Longleftrightarrow\quad
        \Phi_{\mathcal F}(r)=\Phi_{\mathcal F}(r').
\]
A regular local response quotient near $r_0$ is a triple
\[
        U\xrightarrow{\pi}\Qresp_{\Pi,\mathcal F}
        \xrightarrow{\bar\Phi_{\mathcal F}}\Pproc
\]
such that $\pi$ is a surjective submersion, its fibers are exactly the local response fibers, $\Phi_{\mathcal F}|_U=\bar\Phi_{\mathcal F}\circ\pi$, and $\bar\Phi_{\mathcal F}$ is an immersion.
\end{definition}

\begin{proposition}[Local quotient from constant rank]
\label{prop:constant-rank-response-quotient}
Suppose $\Phi_{\mathcal F}$ takes values in a finite-dimensional smooth process chart, is $C^1$, and has constant rank $k$ on a neighborhood of $r_0$.  Then after shrinking the neighborhood it admits a regular local response quotient of dimension $k$.  If the rank changes, no smooth global quotient is implied; the natural response geometry may instead be stratified or singular.
\end{proposition}

\begin{proof}
The constant-rank theorem gives local coordinates $(u,w)$ on the response domain and $y$ on process space in which $\Phi_{\mathcal F}(u,w)=(u,0)$.  Projection $(u,w)\mapsto u$ supplies the local quotient, and $u\mapsto(u,0)$ is the induced immersion.  At a rank-changing point this normal form cannot be chosen with fixed quotient dimension, so smooth-manifold conclusions require an additional stratification analysis.
\end{proof}

\begin{definition}[Pullback process gauge]
Let the represented process displacement space carry a norm $\|\cdot\|_{\rm proc}$.  On a regular quotient define
\[
        \boxed{
        \Fresp(q,v)
        :=
        \left\|D\bar\Phi_{\mathcal F,q}(v)\right\|_{\rm proc}
        }
        \qquad
        (q\in\Qresp_{\Pi,\mathcal F},\ v\in T_q\Qresp_{\Pi,\mathcal F}).
\]
For the finite channel model of Paper~2,
\[
        \|\Delta\|_{\rm proc}:=\frac12\|\Delta\|_\diamond.
\]
\end{definition}

\begin{theorem}[Intrinsic pullback-gauge theorem]
\label{thm:pullback-gauge}
On every regular quotient chart, $\Fresp(q,\cdot)$ is a norm on the tangent space.  It is invariant under smooth reparametrization of quotient coordinates.  If the ambient process norm is not smooth or strictly convex, the result is a possibly nonsmooth Finsler-type structure rather than a classical strongly convex Finsler metric.
\end{theorem}

\begin{proof}
Positive homogeneity and the triangle inequality follow from the process norm and linearity of $D\bar\Phi_q$.  If $\Fresp(q,v)=0$, then $D\bar\Phi_qv=0$; immersion of $\bar\Phi$ gives $v=0$.  Under a quotient-coordinate change $\psi$, the chain rule sends $v$ to $D\psi(v)$ and leaves the process tangent vector $D\bar\Phi(v)$ unchanged, proving invariance.
\end{proof}

\begin{definition}[Intrinsic response length and distance]
For an absolutely continuous quotient path $\gamma:[0,1]\to\Qresp_{\Pi,\mathcal F}$, define
\[
        L_{\rm resp}[\gamma]
        :=\int_0^1\Fresp(\gamma(t),\dot\gamma(t))\,dt,
\]
and
\[
        \dresp(q_0,q_1)
        :=\inf_{\gamma:q_0\to q_1}L_{\rm resp}[\gamma],
\]
with value $+\infty$ if no admissible path connects the endpoints.
\end{definition}

\begin{theorem}[Endpoint process displacement is bounded by response length]
\label{thm:endpoint-process-bound}
For every absolutely continuous response path,
\[
        \dproc\!\left(\bar\Phi(q_0),\bar\Phi(q_1)\right)
        \le L_{\rm resp}[\gamma],
\]
and therefore
\[
        \dproc\!\left(\bar\Phi(q_0),\bar\Phi(q_1)\right)
        \le \dresp(q_0,q_1).
\]
Here $\dproc(\mathcal N_0,\mathcal N_1)=\|\mathcal N_1-\mathcal N_0\|_{\rm proc}$ in an affine process chart.
\end{theorem}

\begin{proof}
The image path $\bar\Phi\circ\gamma$ is absolutely continuous and has speed $\|D\bar\Phi_{\gamma(t)}\dot\gamma(t)\|_{\rm proc}=\Fresp(\gamma(t),\dot\gamma(t))$.  The normed-space triangle inequality bounds endpoint displacement by the integrated speed.  Taking the path infimum gives the second inequality.
\end{proof}

\begin{theorem}[Finite linear chart and Paper~2 recovery]
\label{thm:linear-flat-response-geometry}
Suppose
\[
        \Phi(\mathcal F,r)=\mathcal F+T_\Pi r
\]
on a finite-dimensional affine process chart.  Put
\[
        \bar R_\Pi:=R_\Pi/\ker T_\Pi,
        \qquad
        \bar T_\Pi:\bar R_\Pi\to\operatorname{im}T_\Pi.
\]
Then the quotient response geometry is flat and translation invariant:
\[
        \Fresp([r],[v])=\|\bar T_\Pi[v]\|_{\rm proc},
\]
and straight quotient segments minimize length, so
\[
        \boxed{
        \dresp([r_0],[r_1])
        =\left\|\bar T_\Pi([r_1-r_0])\right\|_{\rm proc}
        }.
\]
If $\|\cdot\|_{\bar c}$ is a calibrated quotient ledger norm, define
\[
L_\Pi^{\rm opt}:=\max_{\|x\|_{\bar c}=1}\|\bar T_\Pi x\|_{\rm proc},
\qquad
\mu_\Pi:=\min_{\|x\|_{\bar c}=1}\|\bar T_\Pi x\|_{\rm proc}.
\]
For a nontrivial quotient,
\[
0<\mu_\Pi\le L_\Pi^{\rm opt}<\infty,
\]
and
\[
\frac{1}{L_\Pi^{\rm opt}}\dresp([r_0],[r_1])
\le
\|[r_1-r_0]\|_{\bar c}
\le
\frac{1}{\mu_\Pi}\dresp([r_0],[r_1]).
\]
\end{theorem}

\begin{proof}
The quotient removes exactly the kernel of $T_\Pi$, so $\bar T_\Pi$ is injective and the displayed tangent gauge follows from \Cref{thm:pullback-gauge}.  The image of a straight quotient segment is a straight process segment, whose normed-space length equals its endpoint displacement.  \Cref{thm:endpoint-process-bound} proves minimality.  Compactness of the quotient unit sphere gives the extremal gains, and injectivity makes the minimum positive.  Homogeneity gives the final comparison.  This is the intrinsic-geometric form of the finite linear response theorem in Paper~2.
\end{proof}

\begin{definition}[Extrinsic free-process monotone]
For a closed free-process set $\mathfrak F_{\rm sep}\subset\Pproc$, define
\[
        M_{\rm sep}(\mathcal N)
        :=\inf_{\mathcal F\in\mathfrak F_{\rm sep}}
        \dproc(\mathcal N,\mathcal F).
\]
\end{definition}

\begin{proposition}[Extrinsic monotone versus intrinsic response distance]
\label{prop:extrinsic-intrinsic}
If $\bar\Phi(q_0)\in\mathfrak F_{\rm sep}$, then
\[
        M_{\rm sep}(\bar\Phi(q_1))\le\dresp(q_0,q_1).
\]
A reverse bound is not automatic.  It requires, for example, that a nearest free process lie in the response chart together with a uniform lower response gain and sufficient quasiconvexity of the response image.
\end{proposition}

\begin{proof}
The free-set infimum is no larger than the distance to the particular free point $\bar\Phi(q_0)$, and \Cref{thm:endpoint-process-bound} bounds that process distance by $\dresp$.  If the nearest free process lies outside the response image, or if the image has long detours relative to ambient process distance, no converse follows.
\end{proof}

\section{Physical-ledger calibration: distance, energy, and the homogeneity fork}
\label{sec:ledger-calibration}

\begin{definition}[Physical interaction-ledger rate]
A physical path ledger on a response quotient is a nonnegative tangent functional $j_q(v)$ whose path cost is
\[
        J[\gamma]=\int_0^1 j_{\gamma(t)}(\dot\gamma(t))\,dt.
\]
The induced directed physical path cost is
\[
        I_\Pi^{\rm path}(q_0,q_1)
        :=\inf_{\gamma:q_0\to q_1}J[\gamma].
\]
The word ``physical'' refers to a separately calibrated resource valuation, not merely to the represented process norm.
\end{definition}

\begin{theorem}[Calibrated ledger realization]
\label{thm:ledger-calibration}
Suppose constants $0<a\le b<\infty$ satisfy
\[
        a\,\Fresp(q,v)
        \le j_q(v)
        \le b\,\Fresp(q,v)
\]
for every admissible tangent direction.  Then
\[
        \boxed{
        a\,\dresp(q_0,q_1)
        \le I_\Pi^{\rm path}(q_0,q_1)
        \le b\,\dresp(q_0,q_1)
        }.
\]
If $j=\Fresp$ exactly, then $I_\Pi^{\rm path}=\dresp$.  Equality is therefore a calibration theorem, not a definition of interaction cost.
\end{theorem}

\begin{proof}
Integrate the pointwise bounds along every admissible path and take the infimum.  Exact calibration makes the two path functionals identical.
\end{proof}

\begin{warning}[When no path metric exists]
Fixed setup burdens, packet thresholds, indivisible tokens, saturation, and history-dependent hysteresis can make physical cost fail local one-homogeneity or refinement additivity.  In those branches $I_\Pi^{\rm path}$ need not be generated by any tangent length.  The response distance remains a valid model geometry, but its identification with the physical ledger fails.
\end{warning}

\begin{theorem}[Local residual energy is squared distance]
\label{thm:residual-energy-distance}
Let $q_0$ be a reference response class and let $\mathcal R$ be a $C^2$ nonnegative residual penalty with
\[
        \mathcal R(q_0)=0,
        \qquad
        D\mathcal R(q_0)=0.
\]
Assume the Hessian
\[
        g^{\rm E}_{q_0}:=D^2\mathcal R(q_0)
\]
is positive definite and extends to a smooth local Riemannian metric $g^{\rm E}$.  Then
\[
        \boxed{
        \mathcal R(q)
        =\frac12 d_{g^{\rm E}}(q_0,q)^2
        +o\!\left(d_{g^{\rm E}}(q_0,q)^2\right)
        }
\]
as $q\to q_0$.
\end{theorem}

\begin{proof}
Choose normal coordinates for $g^{\rm E}$ at $q_0$.  Taylor expansion gives $\mathcal R(q)=\tfrac12 g^{\rm E}_{q_0}(\Delta q,\Delta q)+o(\|\Delta q\|^2)$.  In the same coordinates, squared geodesic distance is $g^{\rm E}_{q_0}(\Delta q,\Delta q)+o(\|\Delta q\|^2)$.  Combining the expansions proves the claim.
\end{proof}

\begin{audit}[Interaction cost is not one object]
Paper~2 contains both branches.  Its quotient coordinate/path-speed bounds are one-homogeneous and naturally define length.  Its smooth residual realization gives a two-homogeneous Hessian energy.  The correct local dictionary is therefore either
\[
        I_\Pi^{\rm path}=\dresp
\]
under exact path-ledger calibration, or
\[
        I_\Pi^{\rm res}=\tfrac12 d_{g^{\rm E}}^2+o(d^2)
\]
under smooth residual realization.  Writing $I_\Pi=d_g$ without specifying the branch loses this distinction.
\end{audit}

\section{Anchoring response geometry to the comparison substrate}
\label{sec:response-anchor}

Response geometry lives on implementation classes.  To obtain substrate geometry one needs a physical map from comparison directions to response directions.

\begin{definition}[Universal infinitesimal comparison family]
Fix a regular region $U\subseteq\Mgeom$.  For each $x\in U$, let
$\mathcal A_x$ denote the full declared family of admitted infinitesimal comparison
continuations based at $x$.  Every $a\in\mathcal A_x$ has a linearized response map
\[
        \ell_{a,x}:T_x\Mgeom\longrightarrow Y_{a,x},
\]
where $Y_{a,x}$ is the corresponding infinitesimal outcome or process-response
space.  Define the universal comparison evaluation map
\[
        E_x:T_x\Mgeom\longrightarrow
        \prod_{a\in\mathcal A_x}Y_{a,x},
        \qquad
        E_xu:=\bigl(\ell_{a,x}u\bigr)_{a\in\mathcal A_x}.
\]
No topology on the full product is needed for the kernel statements below.
On a finite-rank regular branch one may replace the full family by any finite
separating subfamily.
\end{definition}

\begin{definition}[Contextual null distribution]
Two tangent directions are contextually equivalent exactly when every admitted
infinitesimal comparison gives the same response:
\[
        u\approx_x v
        \quad\Longleftrightarrow\quad
        \ell_{a,x}(u-v)=0
        \quad\text{for every }a\in\mathcal A_x.
\]
Therefore the contextual null space is not an independent declaration but the
kernel of the universal comparison evaluation:
\[
        \boxed{
        N_x^{\rm ctx}:=\ker E_x
        =\bigcap_{a\in\mathcal A_x}\ker\ell_{a,x}.}
\]
The contextual tangent quotient is
\[
        T_x^{\rm ctx}\Mgeom:=T_x\Mgeom/N_x^{\rm ctx}.
\]
When $x\mapsto N_x^{\rm ctx}$ has locally constant rank, these fibers form a
smooth quotient bundle.
\end{definition}


\begin{definition}[Comparison and physical-response germs]
\label{def:comparison-response-germs}
Let $\mathcal C_x$ be the set of admissible $C^1$ curve germs
$[\gamma]_x$ through $x$.  Each admitted comparison $a\in\mathcal A_x$ assigns
an outcome germ
\[
        \mathcal L_{a,x}:\mathcal C_x\longrightarrow\mathcal Y_{a,x}.
\]
The complete comparison-germ evaluation is
\[
        \mathcal E_x([\gamma]_x)
        :=\bigl(\mathcal L_{a,x}([\gamma]_x)\bigr)_{a\in\mathcal A_x}.
\]
Let
\[
        \mathcal B_x^{\rm phys}:\mathcal C_x\longrightarrow\mathcal P_x^{\rm red}
\]
be the physical process-response germ valued in the \emph{reduced} response
carrier, obtained after quotienting gauge, auxiliary, and otherwise
record-invisible representation variables.  Thus $\mathcal P_x^{\rm red}$
contains only distinctions that are carried by some physical record sector.
\end{definition}

\begin{theorem}[Closed-world germ-completeness theorem]
\label{thm:closed-world-germ-completeness}
Assume the closed-world closure imported from Papers~0 and~1:
\begin{enumerate}[label=(\roman*)]
\item every finite enforceable distinction belongs to the closed interface ledger
and is therefore readable in the complete reduced process response;
\item no stable reduced process direction is absent from every record sector.
\end{enumerate}
Then the complete comparison-germ evaluation and the reduced physical-response
germ have the same fibers:
\[
\boxed{
\mathcal B_x^{\rm phys}([\gamma]_x)
 =\mathcal B_x^{\rm phys}([\eta]_x)
\iff
\mathcal E_x([\gamma]_x)=\mathcal E_x([\eta]_x).}
\]
Consequently both maps factor through one and the same contextual quotient
\[
        \mathcal Q_x
        :=\mathcal C_x/\!\sim_{\rm ctx},
\qquad
[\gamma]_x\sim_{\rm ctx}[\eta]_x
\iff
\mathcal E_x([\gamma]_x)=\mathcal E_x([\eta]_x),
\]
and the induced reduced physical carrier is canonically isomorphic to
$\mathcal Q_x$ by the unique class-preserving bijection.
\end{theorem}

\begin{proof}
Suppose first that the reduced physical-response germs agree.  Any admitted
comparison distinguishes a finite cost-bearing partition and therefore belongs
to the closed interface ledger.  By closed-world closure its outcome is a
readout of the complete reduced process response.  Equal reduced response germs
therefore give equal outcome germs for every $a$, so the comparison evaluations
agree.

Conversely suppose the comparison evaluations agree but the reduced physical
response germs differ.  Their difference is then a stable distinction in the
reduced process carrier.  Because this carrier has already been quotiented by
record-invisible redundancy, the distinction must be carried by at least one
record sector.  The corresponding finite continuation is an admitted
comparison in the complete family $\mathcal A_x$, and it would distinguish the
two comparison evaluations, contradiction.  Thus the fibers coincide.

Maps with the same fibers induce bijections from their common quotient onto
their images.  Since $\mathcal P_x^{\rm red}$ is defined as the complete reduced
image, the class-preserving bijection is onto and unique.
\end{proof}

\begin{audit}[What closed-world completeness contributes]
The theorem does not infer derivative faithfulness directly from ledger
closure.  It first proves the stronger and correctly typed statement at the
level at which Paper~0 defines physical identity: complete families of finite
admissible continuations.  This avoids the false inference that global
separation automatically implies injective differentiation.  The latter fails,
for example, for $t\mapsto t^2$ at $t=0$.
\end{audit}



\begin{definition}[Finite-stage record test]
\label{def:finite-stage-record-test}
Let $W\subseteq P^{\rm red}$ be a comparison sector, with no topology assumed.
A \emph{finite-stage record test} on $W$ is a map
\[
        r_k:W\longrightarrow S_k,
\]
where $S_k$ is a finite outcome set and $r_k$ is produced by a finite
admissible continuation and deposited in a physical record sector.  The fibres
of all such tests form a Boolean algebra of finite operational distinctions.
They are not, merely by being executable, declared open sets.
\end{definition}

\begin{definition}[Precision-refinement channel and asymptotic quotient]
\label{def:precision-refinement-channel}
A \emph{precision-refinement channel} is a compatible tower
\[
\mathbf r=(r_k,\pi_{k+1,k})_{k\ge0},\qquad
r_k:W\to S_k,\qquad
r_k=\pi_{k+1,k}\circ r_{k+1},
\]
with every \(S_k\) finite. A real realization consists of encodings
\(\iota_k:S_k\to\mathbb R\) for which \(\iota_k r_k\) is locally uniformly
Cauchy with certified error \(\epsilon_k\downarrow0\). Its limit is
\[
\widehat r=\lim_{k\to\infty}\iota_k r_k.
\]
Two complete tower histories are \emph{asymptotically record-equivalent} when
their encoded separation tends to zero. The realization is
\emph{quotient-faithful} when equality of real limits is exactly asymptotic
record-equivalence. Thus finite execution is retained at every precision while
histories that differ only by vanishing boundary conventions are identified in
the reduced channel.
\end{definition}

\begin{proposition}[Unreduced tower-faithfulness no-go]
\label{prop:tower-faithful-no-go}
Let \(X_{\mathbf r}\) be the inverse limit of the finite discrete stage spaces.
If a real realization is continuous and injective on the unreduced tower space
\(X_{\mathbf r}\), then its image is compact and zero-dimensional. A finite
product of such tower-faithful channel images is again zero-dimensional and
cannot contain a nontrivial connected smooth continuation orbit.
\end{proposition}

\begin{proof}
A countable inverse limit of finite discrete spaces is compact, metrizable, and
zero-dimensional. A continuous injection from a compact space into Hausdorff
\(\mathbb R\) is a homeomorphism onto its image, so the image remains
zero-dimensional. Finite products preserve zero-dimensionality. Every connected
subset is therefore a singleton.
\end{proof}

\begin{corollary}[Positive-dimensional smooth continuum requires more than faithful tower relabelling]
\label{cor:continuum-needs-quotient}
A continuous injective real relabelling of an unreduced profinite tower may
exist, but its image remains zero-dimensional.  Therefore no finite product of
such faithful relabellings can by itself contain a nontrivial connected
positive-dimensional smooth continuation orbit.  The physical construction
must either pass to an asymptotic quotient that changes the relevant topology
or add independently certified non-profinite smooth structure.  Real-valued
faithful embeddings of Cantor-type tower spaces are not excluded.
\end{corollary}

\begin{definition}[Record-limit topology]
\label{def:record-limit-topology}
Let $\mathscr R_{\rm ch}(W)$ be the family of ledger-certified real channel
limits on $W$.  The \emph{record-limit topology} $\tau_{\rm rec}$ is the
initial topology for the maps
\[
        \widehat r:W\longrightarrow\mathbb R,
        \qquad \widehat r\in\mathscr R_{\rm ch}(W),
\]
where $\mathbb R$ carries its ordinary topology.  Equivalently, a subbasis is
formed by
\[
        \widehat r^{-1}((a,b)),
        \qquad a,b\in\mathbb Q,
        \quad a<b.
\]
Finite-stage fibres certify finite-resolution distinctions and approximate
these neighbourhoods; they are not automatically clopen in $\tau_{\rm rec}$.
\end{definition}

\begin{warning}[Why the finite-stage topology is rejected]
If every fibre of every finite-stage test were declared open, the resulting
Hausdorff refinement topology would be zero-dimensional and totally
disconnected.  It could not contain a nontrivial connected smooth continuation
orbit.  APF therefore uses the finite stages as executable approximants and
constructs physical neighbourhoods from the ordered limits of precision
channels.  This is the exact place where the continuum topology enters; it is
not the unrestricted profinite completion of the finite partitions.
\end{warning}

\begin{proposition}[Finite-stage clopen no-go]
\label{prop:finite-stage-clopen-no-go}
Let $\tau_{\rm fin}$ be the topology generated by declaring every finite-stage
record fibre open.  Then every such fibre is clopen.  If the finite stages
separate points, $(W,\tau_{\rm fin})$ is Hausdorff and zero-dimensional; hence
all connected subsets are singletons.  No positive-dimensional connected
smooth orbit can be embedded continuously in this topology.
\end{proposition}

\begin{proof}
Each stage has finitely many fibres, so the complement of any one fibre is the
finite union of the others and is open.  Thus the fibre basis is clopen and the
space is zero-dimensional.  A Hausdorff zero-dimensional space is totally
disconnected.  The continuous image of a connected manifold is connected, so
an injective continuous orbit inclusion would have singleton image unless the
orbit were zero-dimensional.
\end{proof}

\begin{proposition}[Closed-world Hausdorff separation]
\label{prop:closed-world-hausdorff}
If the complete real channel family separates reduced physical states, then
$(P^{\rm red},\tau_{\rm rec})$ is Hausdorff.
\end{proposition}

\begin{proof}
Let $p\neq q$.  Closed-world reduction supplies a real precision channel
$\widehat r$ with $\widehat r(p)\neq\widehat r(q)$.  Choose disjoint rational
open intervals containing the two values.  Their inverse images are disjoint
$\tau_{\rm rec}$-neighbourhoods of $p$ and $q$.
\end{proof}

\begin{assumption}[Countable channel-specification grammar]
\label{ass:countable-channel-grammar}
The primitive instruction tokens, finitely encoded exact parameters, stopping
rules, bonding-map descriptions, encodings, and error certificates available
on the declared interface form a countable specification language.  Arbitrary
real constants are not admitted as single primitive tokens without an explicit
countable code.
\end{assumption}

\begin{definition}[Finitely specified physical channel]
\label{def:countable-refinement-protocol}
A physical precision channel is \emph{finitely specified} when one finite word
in the grammar of Assumption~\ref{ass:countable-channel-grammar} specifies its
stage tests, bonding maps, requested-precision stopping rule, encodings, and
certified error bound.
\end{definition}

\begin{theorem}[Countable specified channel family]
\label{thm:countable-channel-family}
Under Assumption~\ref{ass:countable-channel-grammar}, the finitely specified
physical precision channels on a fixed closed interface form a countable family.
\end{theorem}
\begin{proof}
Finite words over a countable alphabet form a countable set.
\end{proof}

\begin{gate}[Countable topology generation]
\label{gate:countable-topology-generation}
The finitely specified quotient-faithful real limits contain a countable
subfamily whose initial topology equals the record-limit topology of the
complete ledger-certified channel family.
\end{gate}

\begin{corollary}[Conditional second countability]
\label{prop:record-second-countable}
Under Assumption~\ref{ass:countable-channel-grammar} and
Gate~\ref{gate:countable-topology-generation}, the record-limit topology is
second countable.
\end{corollary}
\begin{proof}
Inverse images of rational intervals under the countable topology-generating
subfamily form a countable subbasis, and finite intersections form a countable
base.
\end{proof}

\begin{lemma}[Finite stages approximate record-limit neighbourhoods]
\label{lem:uniform-limits-no-refinement}
Let $\widehat r$ be a channel limit with certified local error
$|\iota_k r_k-\widehat r|\le\epsilon_k$.  For every $p$ and every open interval
$I\ni\widehat r(p)$, there are a smaller interval $J$ with
$\widehat r(p)\in J\Subset I$ and a stage $k$ such that the finite-stage cell
containing $p$ is mapped by $\iota_k$ into the $\epsilon_k$-thickening of $J$,
which lies in $I$.  Hence every record-limit neighbourhood is operationally
resolvable to finite certified precision.
\end{lemma}

\begin{proof}
Choose $J$ with positive distance $\delta$ from its closure to
$\mathbb R\setminus I$.  Take $k$ with $\epsilon_k<\delta/2$ and with
$\iota_k r_k(p)$ within $\epsilon_k$ of $\widehat r(p)$.  Compatibility and
the certified error bound place all states in the corresponding stage cell
whose encoded value lies in the chosen thickening inside $I$.  No claim that
the entire stage cell is open is required.
\end{proof}

\begin{definition}[Executable channel algebra]
\label{def:finite-record-algebra-topology}
For $U\in\tau_{\rm rec}$, let $\mathscr A_{\rm ch}(U)$ be the unital real
algebra generated by restrictions of real precision-channel limits and by
ledger-certified locally uniform limits of executable scalar postprocessings.
Every generator is $\tau_{\rm rec}$-continuous by construction or by uniform
convergence.  Let $\mathscr A_0(U)$ denote its closure under such locally
uniform limits.
\end{definition}

\begin{audit}[Topology ledger]
The finite stages provide the finite executable evidence.  The real channel
limits provide the ordered topology.  Closed-world completeness provides point
separation.  A countable specification grammar provides a countable specified channel family; second
countability follows only when that family is topology-generating.  None of these statements identifies the raw inverse limit of
finite partitions with a continuum, and none assumes a manifold.
\end{audit}

\begin{definition}[Executable postprocessing and smooth saturation]
\label{def:executable-vs-saturated}
Let $f=(f_1,\ldots,f_m)$ be a jointly realizable finite channel map at the requested precision.  A scalar
postprocessing is \emph{executable} when it is given by a finite protocol at
each requested precision.  It is \emph{observationally conservative} when it
factors through $f$, and therefore cannot distinguish two states that $f$
does not already distinguish.

The \emph{smooth saturation} of a family of real readouts is its closure under
all maps $\psi\in C^\infty(\mathbb R^m)$.  Smooth saturation is a mathematical
closure of the represented observable algebra; it does not assert that every
smooth map is itself a primitive finite program.
\end{definition}

\begin{theorem}[Lossless postprocessing factor theorem]
\label{thm:lossless-postprocessing-factor}
Every executable deterministic postprocessing of a closed joint record is
observationally conservative.  Conversely, adjoining any scalar function already constant on the fibres of the closed joint record changes neither those fibres nor the closed-world contextual quotient.  No executability claim is implied in this converse.
\end{theorem}

\begin{proof}
If $h=\varphi\circ f$, then $f(p)=f(q)$ implies $h(p)=h(q)$, so $h$ creates no
new distinction class.  Appending an executable protocol to the record
continuation changes implementation cost but not the partition of physical
states induced by the complete record.  Conversely, any scalar function
constant on every fibre of $f$ factors set-theoretically through the image of
$f$; adjoining it leaves those fibres unchanged.  Hence the contextual
quotient is invariant under such lossless relabelling.
\end{proof}

\begin{definition}[Refinement-complete APF differential sheaf]
\label{def:record-sheaf}
For $U\in\tau_{\rm rec}$, let $\mathscr R(U)$ contain the refinement-stable
real readouts obtained from executable records and their ledger-certified
local limits.  Define $\mathscr C_{\rm APF}(U)$ to be the localization and
smooth saturation of $\mathscr R(U)$: it is the smallest family of real
functions on $U$ that
\begin{enumerate}[label=(\roman*)]
\item contains $\mathscr R(U)$;
\item is closed under restriction and local gluing;
\item satisfies
\[
f_1,\ldots,f_m\in\mathscr C_{\rm APF}(U),\quad
\psi\in C^\infty(\mathbb R^m)
\Longrightarrow
\psi(f_1,\ldots,f_m)\in\mathscr C_{\rm APF}(U).
\]
\end{enumerate}
\end{definition}

\begin{theorem}[Differential-space closure theorem]
\label{thm:apf-cinfty-closure}
The triple
\[
(P^{\rm red},\tau_{\rm rec},\mathscr C_{\rm APF})
\]
is a Sikorski differential space.  Its topology is the initial topology of the
physical readouts, and its differential structure is closed under smooth
superposition and localization.
\end{theorem}

\begin{proof}
The record-limit topology is, by Definition~\ref{def:record-limit-topology},
the initial topology of the generating physical channel limits.  Smooth
superpositions of continuous generators remain continuous, and localization
adds only functions that locally agree with such continuous functions.
Definition~\ref{def:record-sheaf} is therefore closed under smooth
superposition and localization and induces the same topology.  These are the
Sikorski axioms.  Proposition~\ref{prop:closed-world-hausdorff} and
Proposition~\ref{prop:record-second-countable} supply Hausdorffness and, on
countably generated sectors, second countability; neither property is needed
to define the differential structure itself.  The lossless-postprocessing
factor theorem shows that the saturation changes no closed-world record fibre.
\end{proof}

\begin{audit}[What is derived and what is saturated]
The executable channel algebra, its record-limit topology, and invariance of
record fibres under deterministic postprocessing are derived from APF closure.
Closure under the full $C^\infty$ functional calculus is the canonical smooth
saturation of that algebra.  It is not disguised computability and it does not
assert that finite records alone force a manifold.  Its role is to define the
minimal differential structure compatible with the already represented
readouts; manifold structure appears only on the regular continuation orbits
below.
\end{audit}

\begin{definition}[Derivation tangent and algebraic cotangent]
\label{def:derivation-tangent}
Let $\mathscr C_p$ be the stalk of $\mathscr C_{\rm APF}$ at $p$, and let
$\mathfrak m_p=\{[f]_p:f(p)=0\}$ be its maximal ideal.  The primary tangent is
\[
        T_p^{\rm der}P
        :=\operatorname{Der}_p(\mathscr C_p,\mathbb R),
\]
the real derivations satisfying Leibniz' rule.  The algebraic first-order
cotangent candidate is
\[
        T^{*,\rm alg}_pP:=\mathfrak m_p/\mathfrak m_p^2.
\]
Every derivation defines a linear functional on
$\mathfrak m_p/\mathfrak m_p^2$; Proposition~\ref{prop:derivation-cotangent-pairing}
shows that this pointwise pairing is canonical and perfect.  Singular behaviour
appears through rank variation and failure of finite local presentation.
\end{definition}

\begin{proposition}[Canonical tangent--cotangent duality]
\label{prop:derivation-cotangent-pairing}
Evaluation on first-order classes gives a canonical vector-space isomorphism
\[
 T_p^{\rm der}P
 \cong
 \bigl(\mathfrak m_p/\mathfrak m_p^2\bigr)^*,
 \qquad
 X\longmapsto\bigl([f]\mapsto Xf\bigr).
\]
Singularity is therefore detected by failure of finite local presentation,
rank jumps, or variation of this common dimension--not by a spurious gap
between the two objects.
\end{proposition}

\begin{proof}
Leibniz' rule makes every derivation vanish on $\mathfrak m_p^2$, so the map is
well defined and injective.  Conversely, for
$\lambda\in(\mathfrak m_p/\mathfrak m_p^2)^*$ define
\[
X_\lambda(f):=\lambda\bigl([f-f(p)]\bigr).
\]
Modulo $\mathfrak m_p^2$ one has
$fg-f(p)g(p)=f(p)(g-g(p))+g(p)(f-f(p))$, hence $X_\lambda$ satisfies Leibniz'
rule at $p$.  The two constructions are inverse.
\end{proof}

\begin{definition}[Finite differential presentation]
\label{def:locally-finitely-generated}
A record-open neighbourhood $U$ has a \emph{finite differential presentation}
when there exist real channel limits $f_1,\ldots,f_N\in\mathscr C_{\rm APF}(U)$ such that
\begin{enumerate}[label=(\roman*)]
\item the evaluation map
\[
F_U=(f_1,\ldots,f_N):U\to\mathbb R^N
\]
is injective and is a homeomorphism onto its image;
\item every $g\in\mathscr C_{\rm APF}(U)$ is locally of the form
$g=\psi\circ F_U$ for some $\psi\in C^\infty(\mathbb R^N)$.
\end{enumerate}
The locus of points admitting such a neighbourhood is denoted
$P^{\rm fdiff}$.
\end{definition}

\begin{definition}[Ledger-independent standing channels]
\label{def:independent-standing-channel}
A jointly realizable family of precision channels
$\mathbf r^1,\ldots,\mathbf r^m$ is \emph{ledger-independent} on $U$ when no
complete tower $\mathbf r^i$ is determined by the joint complete tower of the
others.  Equivalently, for every $i$ there exist reduced states $p_i,q_i\in U$
that agree in all channels $j\neq i$ but differ in channel $i$.  A
ledger-independent channel therefore carries a distinction not recoverable
from the remaining standing commitments.
\end{definition}

\begin{theorem}[Additive channel-capacity bound]
\label{thm:additive-channel-capacity}
Assume the APF positive standing-cost floor $\varepsilon_*>0$, additive ledger
accounting for jointly realized ledger-independent commitments, and finite
local capacity $C(U)<\infty$.  Then every ledger-independent jointly realizable
family of precision channels on $U$ has cardinality
\[
        m\le \left\lfloor\frac{C(U)}{\varepsilon_*}\right\rfloor.
\]
The bound concerns the number of independent standing channels, not the number
of finite-stage outcomes, the number of states, or the cardinality of any
refinement limit.
\end{theorem}

\begin{proof}
Ledger independence assigns to each channel a distinction unavailable from the
others.  By the positive floor, maintaining each such standing distinction
costs at least $\varepsilon_*$.  Additivity for simultaneously maintained
independent commitments therefore gives total standing charge at least
$m\varepsilon_*$.  Admissibility requires $m\varepsilon_*\le C(U)$.
\end{proof}

\begin{definition}[Complete channel basis]
\label{def:complete-channel-basis}
A \emph{complete channel basis} on $U$ is a jointly realizable channel family
that separates reduced states and is inclusion-minimal with that property.
Equivalently, deleting any one channel destroys separation.
\end{definition}

\begin{lemma}[Minimal completeness implies ledger independence]
\label{lem:minimal-complete-independent}
Every complete channel basis is ledger-independent.
\end{lemma}

\begin{proof}
If channel $i$ were determined by all the others, then equality in the
remaining channels would force equality in channel $i$.  Deleting channel $i$
would therefore preserve the joint fibres and hence preserve separation,
contradicting inclusion-minimality.
\end{proof}

\begin{theorem}[Finite-channel completeness]
\label{thm:finite-channel-completeness}
Let $U$ admit a complete channel basis and satisfy the additive
channel-capacity theorem.  Then every complete channel basis is finite:
\[
        \mathbf F_U=(\mathbf r^1,\ldots,\mathbf r^N),
        \qquad
        N\le\left\lfloor\frac{C(U)}{\varepsilon_*}\right\rfloor.
\]
Its joint real limit map
\[
        \widehat F_U=(\widehat r^1,\ldots,\widehat r^N):U\to\mathbb R^N
\]
is injective whenever the chosen channels possess faithful real realizations.
\end{theorem}

\begin{proof}
By Lemma~\ref{lem:minimal-complete-independent}, a complete channel basis is
ledger-independent.  The additive capacity bound therefore makes it finite.
Completeness says that equality of every complete channel tower implies equality
in the reduced carrier.  Faithful real realization preserves each tower's complete fibres; hence equality of all real limits implies equality of all channel towers and therefore equality of states.  Thus $\widehat F_U$ is injective.
\end{proof}

\begin{gate}[Finite closed-interface generation]
\label{gate:finite-closed-interface-generation}
Finite capacity bounds every jointly realizable ledger-independent channel
family, but it does not by itself prove that one finite family separates the
entire reduced sector. The missing statement is the finite-generation theorem
\[
\boxed{\text{closed-world completeness + finite capacity}
\Longrightarrow\text{finite complete physical channel basis}.}
\]
It requires either a channel-dependence exchange theorem, a compact finite-subcover
theorem for physical distinctions, or an equivalent APF generation result.
Without one of these, a countable separating family can require infinitely many
channels even though no large independent family is jointly realizable.
\end{gate}

\begin{proposition}[Finite-channel sufficiency versus finite differential sufficiency]
\label{prop:finite-sufficiency-split}
Theorem~\ref{thm:finite-channel-completeness} gives a finite separating channel
map $\widehat F_U$.  It gives a finite differential presentation exactly when
\begin{enumerate}[label=(\roman*)]
\item $\widehat F_U$ is a homeomorphism from $U$ onto its image in the record
topology; and
\item every $g\in\mathscr C_{\rm APF}(U)$ factors locally as
$g=\psi\circ\widehat F_U$ with $\psi\in C^\infty(\mathbb R^N)$.
\end{enumerate}
\end{proposition}

\begin{proof}
Separation gives injectivity.  The topology and smooth-factorization clauses
are precisely the two additional conditions in
Definition~\ref{def:locally-finitely-generated}.  No finite-image inference is
used.
\end{proof}

\begin{theorem}[Subcartesian realization criterion]
\label{cor:subcartesian-realization}
Every point of $P^{\rm fdiff}$ has a neighbourhood isomorphic, as a
differential space, to a differential subspace of a subset of
$\mathbb R^N$.  Hence $P^{\rm fdiff}$ is locally subcartesian.  Conversely,
any locally subcartesian presentation by physical readouts is a finite
differential presentation in the sense above.
\end{theorem}

\begin{proof}
Condition (i) identifies $U$ homeomorphically with $F_U(U)$.  Condition (ii)
says that the differential structure on $U$ is exactly the pullback of the
restriction of $C^\infty(\mathbb R^N)$ to that image.  This is the defining
local model of a subcartesian differential space.  The converse is immediate
from a finite ambient-coordinate presentation.
\end{proof}

\begin{gate}[Finite smooth generation]
\label{gate:finite-differential-sufficiency}
Even after a finite state-separating channel family is obtained, finite smooth
generation is a distinct theorem. It requires that the continuation-regular
readout sheaf be locally the smooth functional calculus of that finite map:
\[
\mathscr C_{\rm APF}|_U
=\operatorname{Loc}\,C^\infty(f_1,\ldots,f_N).
\]
State separation alone does not imply this equality. Raw threshold records may
be discontinuous, and continuous limit readouts need not be smooth functions of
a separating map. The subcartesian orbit theorem is invoked only after this
finite smooth-generation statement closes.
\end{gate}

\begin{audit}[Last-open-item ledger]
The topology layer now has no hidden certificate: quotient-faithful channel
limits, Hausdorff separation, and second countability are derived. Two
substantive APF theorems remain before smooth finite-dimensional geometry is
unconditional:
\begin{enumerate}[label=(\roman*)]
\item finite closed-interface generation of a state-separating physical channel
basis;
\item finite smooth generation of the continuation-regular readout sheaf by
that basis.
\end{enumerate}
These are not notation problems. They are the exact remaining continuum and
regularity content.
\end{audit}

\begin{remark}[Imported singular differential geometry]
For subcartesian differential spaces, derivations satisfy the smooth chain
rule.  The orbits of a family complete in the orbit-theorem sense carry unique
smooth manifold structures for which the orbit inclusions are smooth
immersions.  The paper imports these established results rather than
reproving them.  See Sikorski's differential-space formalism and the subcartesian orbit results of Śniatycki and collaborators.\footnote{J. Śniatycki, \emph{Orbits of families of vector fields on subcartesian spaces}, Ann. Inst. Fourier 53 (2003), 2257--2296, DOI: \href{https://doi.org/10.5802/aif.2006}{10.5802/aif.2006}; J. Śniatycki, \emph{Differential Geometry of Singular Spaces and Reduction of Symmetry}, Cambridge University Press (2013), DOI: \href{https://doi.org/10.1017/CBO9781139136990}{10.1017/CBO9781139136990}; R. Cushman and J. Śniatycki, \emph{On subcartesian spaces Leibniz' rule implies the chain rule}, Can. Math. Bull. 63 (2020), DOI: \href{https://doi.org/10.4153/S0008439519000513}{10.4153/S0008439519000513}.}
\end{remark}

\begin{definition}[Continuation distribution and geometric orbits]
\label{def:continuation-distribution}
Let $\mathscr X_{\rm adm}$ be a family of vector fields on
$P^{\rm fdiff}$, where a vector field is a derivation that generates a local
one-parameter group of local diffeomorphisms.  The family is
\emph{complete in the orbit-theorem sense} when for every
$X,Y\in\mathscr X_{\rm adm}$ and every time $t$ for which the local flow
$\varphi_t^X$ is defined, the pushforward
$(\varphi_t^X)_*Y$ agrees locally with a member of
$\mathscr X_{\rm adm}$.  Equivalently, the family is closed locally under the
flow transports used to form its pseudogroup.  Define
\[
        \Delta_p:=\operatorname{span}\{X(p):X\in\mathscr X_{\rm adm}\}
        \subseteq T_p^{\rm der}P.
\]
The orbit through $p$ under finite compositions of these local flows is denoted
$\mathcal O_p$.
\end{definition}

\begin{theorem}[Smooth geometric-orbit theorem]
\label{thm:smooth-geometric-orbit}
Every orbit $\mathcal O_p$ of an admissible continuation family complete in the orbit-theorem sense on $P^{\rm fdiff}$ carries a unique smooth manifold structure for which its inclusion into the differential space is a smooth immersion.  Its tangent space is the continuation distribution:
\[
        T_q\mathcal O_p=\Delta_q
        \qquad(q\in\mathcal O_p).
\]
Thus smooth manifolds arise as the dynamically connected regular pieces of the
APF differential space, not as prior substrate assumptions.
\end{theorem}

\begin{proof}
By Theorem~\ref{cor:subcartesian-realization}, $P^{\rm fdiff}$ is locally a
subcartesian differential space.  The source-level completeness clause in Definition~\ref{def:continuation-distribution} is exactly the closure hypothesis required by the orbit theorem for complete families of vector fields on subcartesian spaces; hence it applies to $\mathscr X_{\rm adm}$ and
makes each orbit a smooth manifold with the stated tangent distribution.
\end{proof}

\begin{definition}[Maximal smooth geometric branch and singular locus]
\label{def:maximal-smooth-branch}
A \emph{smooth geometric branch} is a maximal connected continuation orbit on
which $\dim\Delta$ is locally constant and positive.  The union of all such
orbits is $P^{\rm geom}_{\rm reg}$.  Its complement
\[
        P^{\rm sing}:=P^{\rm red}\setminus P^{\rm geom}_{\rm reg}
\]
is the intrinsic singular locus.  It includes zero-rank, rank-jump,
saturation, branching, and non-flowable points.
\end{definition}

\begin{theorem}[Existence of the regular geometric regime]
\label{thm:primitive-regular-regime}
Let $p\in P^{\rm fdiff}$.  If the admissible continuation family is complete in the orbit-theorem sense near $p$, continuation-active at $p$, and has locally constant
positive orbit rank, then there is a finite-dimensional smooth geometric
branch through $p$.
\end{theorem}

\begin{proof}
The subcartesian realization criterion supplies the required local singular
smooth model.  Continuation activity gives a nonzero flow-generating vector
field, and Theorem~\ref{thm:smooth-geometric-orbit} gives an immersed smooth
orbit through $p$.  Local constancy of $\dim\Delta$ fixes its dimension, while
maximal connected extension produces the branch.
\end{proof}

\begin{theorem}[Consolidated APF regular-geometry theorem]
\label{thm:consolidated-regular-geometry}
Let $U\subseteq P^{\rm red}$ be a closed-world comparison domain satisfying:
a quotient-faithful state-separating real channel family; the countable grammar
and topology-generation conditions; a complete finite channel basis; finite
smooth generation; and a locally flowable continuation family complete in the
orbit-theorem sense.  Then:
\begin{enumerate}[label=(\roman*)]
\item the channel limits generate a Hausdorff second-countable record topology;
\item localization and smooth saturation produce a Sikorski differential space;
\item on $U\cap P^{\rm fdiff}$ the space is locally subcartesian;
\item every orbit of the complete admissible-continuation family is an immersed
smooth manifold; and
\item on each constant-rank orbit $S$ satisfying
Gate~\ref{gate:smooth-response-factor-completeness},
\[
\ker D\beta_x
=N_x^{\rm ctx}
=\bigcap_{a\in\mathcal A_x}\ker\ell_{a,x},
\qquad x\in S,
\]
with injective induced map
\[
T_xS/N_x^{\rm ctx}\hookrightarrow T_{\beta(x)}\mathcal O_{\beta(x)}.
\]
\end{enumerate}
Thus smooth response geometry is the regular orbit locus of the APF record
differential space, not an upstream manifold postulate.
\end{theorem}

\begin{proof}
Item (i) is Propositions~\ref{prop:closed-world-hausdorff} and
\ref{prop:record-second-countable}.  Items (ii)--(iv) are
Theorem~\ref{thm:apf-cinfty-closure},
Theorem~\ref{cor:subcartesian-realization}, and
Theorem~\ref{thm:smooth-geometric-orbit}.  Item (v) is
Theorem~\ref{thm:regular-stratum-closed-world-kernel} and
Corollary~\ref{cor:faithful-quotient-response}.
\end{proof}

\begin{audit}[Resolution of the continuum-emergence objection]
The raw refinement completion remains profinite and may be totally
disconnected.  Smooth continuum is not obtained by declaring its points
close.  It appears on the finite-differential-presentation locus as an orbit of
physically admissible local continuation derivations.  Thus two independent
conditions are visible rather than hidden: smooth local factorization of the
record algebra, and continuation-active local flows.  Failure of either leaves
a singular, higher-order, or zero-dimensional record sector rather than a
smuggled manifold.
\end{audit}

\begin{audit}[Fortification ledger: what remains conditional]
The differential-space architecture is now noncircular, but four obligations
remain deliberately named rather than hidden:
\begin{enumerate}[label=(\roman*)]
\item the complete physical channel family must admit faithful real
realizations and separate reduced states;
\item the sector must admit a countable topology-generating channel subfamily
when second countability is claimed;
\item finite-channel Regime~R must supply a complete channel basis, while the
capacity theorem proves only that such a basis is finite;
\item the finite joint channel map must satisfy the homeomorphism and smooth
local-factorization clauses before the subcartesian orbit theorem is invoked.
\end{enumerate}
These are falsifiable regime certificates.  None is replaced by the statement
that finite partitions alone produce a continuum.
\end{audit}

\begin{definition}[Contextual comparison differentials]
\label{def:contextual-comparison-differentials}
On a smooth geometric branch $S$, every admitted regularized comparison
$L_a\in\mathscr C_{\rm APF}$ has differential
\[
        \ell_{a,x}:=d(L_a|_S)_x:T_xS\to Y_{a,x}.
\]
Define
\[
        N_x^{\rm ctx}:=\bigcap_{a\in\mathcal A_x}\ker\ell_{a,x}.
\]
Let $\beta:S\to P^{\rm red}$ be the complete reduced physical-response map and
$B_x^{\rm phys}:=D\beta_x$ on the orbit manifold.
\end{definition}

\begin{audit}[Same fibres do not imply the same derivative kernel]
On $\mathbb R$, $L(t)=t$ and $\beta(t)=t^3$ have identical fibres, but
$dL_0\neq0$ and $D\beta_0=0$.  Set-level closed-world completeness is therefore
insufficient for first-order faithfulness.
\end{audit}

\begin{gate}[Smooth response-factor completeness]
\label{gate:smooth-response-factor-completeness}
On a neighbourhood $V\subseteq S$ of $x$, there is a finite APF readout
generator
\[
F=(f_1,\ldots,f_N):V\to\mathbb R^N
\]
whose components are admitted readouts and generate the local Sikorski
structure, together with smooth maps on the relevant images such that
\[
F=\Psi\circ\beta,
\qquad
\beta=\Theta\circ F.
\]
\end{gate}

\begin{theorem}[Regular-orbit closed-world kernel theorem]
\label{thm:regular-stratum-closed-world-kernel}
On every smooth geometric branch satisfying closed-world germ completeness and
Gate~\ref{gate:smooth-response-factor-completeness},
\[
\boxed{
\ker D\beta_x
=\ker DF_x
=\bigcap_{a\in\mathcal A_x}\ker\ell_{a,x}
=N_x^{\rm ctx}.}
\]
\end{theorem}

\begin{proof}
Differentiating $F=\Psi\circ\beta$ and $\beta=\Theta\circ F$ gives
$\ker D\beta_x=\ker DF_x$.  Every admitted readout is locally a smooth
function of $F$, so $\ker DF_x$ lies in every admitted readout kernel.
Conversely the components of $F$ are admitted readouts, so a vector
annihilated by all admitted readout differentials lies in $\ker DF_x$.
\end{proof}

\begin{corollary}[Faithful quotient response]
\label{cor:faithful-quotient-response}
The induced map
\[
 \widetilde B_x:T_xS/N_x^{\rm ctx}
 \longrightarrow T_{\beta(x)}\mathcal O_{\beta(x)},
 \qquad [u]\longmapsto B_x^{\rm phys}u,
\]
is injective.
\end{corollary}

\begin{definition}[Universal operational response carrier]
\label{def:universal-operational-response-carrier}
For $x\in S$, define
\[
        T_x^{\rm op}S:=T_xS/N_x^{\rm ctx}.
\]
This quotient is universal among linear first-order response realizations that
annihilate every contextually null direction.
\end{definition}

\begin{theorem}[Exact-realization criterion for a concrete response model]
\label{thm:exact-realization-criterion}
Let $B_x^{\rm raw}:T_xS\to W_x$ be a concrete finite response model.  The
following are equivalent:
\begin{enumerate}[label=(\roman*)]
\item $\ker B_x^{\rm raw}=N_x^{\rm ctx}$;
\item $B_x^{\rm raw}$ factors through an injective map
$T_x^{\rm op}S\hookrightarrow W_x$;
\item the model is sound on contextual null directions and jointly separating
on the quotient.
\end{enumerate}
\end{theorem}

\begin{proof}
This is the first isomorphism theorem applied to the quotient map
$T_xS\to T_xS/N_x^{\rm ctx}$.
\end{proof}

\begin{theorem}[Canonical response-orbit quotient gauge]
\label{thm:canonical-operational-anchor}
Let $S$ be a smooth geometric branch and suppose the process carrier has norm
$\|\cdot\|_{\rm proc}$.  Then
\[
        c_x^{\rm op}([u])
        :=\|\widetilde B_x[u]\|_{\rm proc}
\]
defines a positive homogeneous gauge on $T_x^{\rm op}S=T_xS/N_x^{\rm ctx}$.  If the response norm
is smooth and strongly convex away from zero, these gauges form a Finsler
structure on each locally constant-rank orbit stratum.  This gauge is intrinsic
to the response orbit; it is not yet a map from physical site displacements.
\end{theorem}

\begin{proof}
Well-definedness follows from quotienting by the kernel; positivity follows
from injectivity of $\widetilde B_x$.  Homogeneity and convexity pull back from
the process norm.  Smooth bundle dependence follows on constant-rank orbit
strata.
\end{proof}

\begin{corollary}[Direct pullback case]
\label{cor:direct-pullback-anchor}
When $\beta$ itself is the complete response map, the operational gauge is the
direct pullback of the process norm by $D\beta$ after quotienting its kernel.
\end{corollary}

\begin{theorem}[Intrinsic response-orbit length theorem]
\label{thm:response-orbit-length}
For every piecewise $C^1$ curve $\delta$ in a constant-rank smooth response
orbit,
\[
 L_{\rm orb}[\delta]
 =\int c^{\rm op}_{\delta(t)}([\dot\delta(t)])\,dt
 =L_{\rm proc}[\beta\circ\delta].
\]
Consequently endpoint process displacement is bounded by intrinsic
response-orbit length, and the induced point function is a pseudometric before
quotienting zero-length connectivity.
\end{theorem}
\begin{proof}
The chain rule gives
$D\beta_{\delta(t)}\dot\delta(t)=(\beta\circ\delta)'(t)$.  Substitute into the
definition of $c^{\rm op}$ and integrate.
\end{proof}

\begin{gate}[Physical comparison-substrate anchor]
\label{gate:physical-substrate-anchor}
There exist a physical regular comparison domain $M_{\rm reg}$, a smooth site
map $s:M_{\rm reg}\to S$, and a locally constant-rank bundle map
\[
A_x:T_xM_{\rm reg}\to T_{s(x)}^{\rm op}S
\]
such that, with $q_y:T_yS\to T_y^{\rm op}S$ the quotient map:
\begin{enumerate}[label=(\roman*)]
\item $\ker A_x$ is exactly the physical comparison-null subspace;
\item $A_x$ has declared process/ledger provenance and is response-gauge
invariant;
\item $A_x=q_{s(x)}\circ Ds_x$;
\item $c_x^{\rm dir}(u)=c_{s(x)}^{\rm op}(A_xu)$; and
\item the calibration is stable on the regular branch.
\end{enumerate}
\end{gate}

\begin{theorem}[Conditional physical substrate pullback]
\label{thm:anchored-cost-geometry}
If Gate~\ref{gate:physical-substrate-anchor} holds, define
$c_x^{\rm sub}(u):=c_{s(x)}^{\rm op}(A_xu)$.  Then every piecewise $C^1$ curve
$\gamma$ in $M_{\rm reg}$ satisfies
\[
L_{\rm sub}[\gamma]
:=\int c^{\rm sub}_{\gamma(t)}(\dot\gamma(t))\,dt
=L_{\rm orb}[s\circ\gamma].
\]
\end{theorem}
\begin{proof}
Gate~\ref{gate:physical-substrate-anchor}(iii) gives
$A_{\gamma(t)}\dot\gamma(t)=[(s\circ\gamma)'(t)]$.  Equality of the integrands,
and hence of the lengths, follows.
\end{proof}

\begin{proposition}[Positive cost metric from quadraticity]
\label{prop:positive-cost-metric}
If the process norm satisfies the parallelogram identity on the image of every
$\widetilde B_x$, then polarization defines a positive-definite inner product
on $T_x^{\rm op}S$.  These inner products form the positive cost metric
$g^{\rm cost}$ on each smooth constant-rank branch.
\end{proposition}

\begin{proof}
The Jordan--von Neumann theorem turns the quotient norm into an inner product.
Injectivity makes it positive definite; smooth dependence follows from the
smooth response representation.
\end{proof}

\begin{definition}[Comparison atlas]
The atlas used below is the maximal smooth atlas of the continuation orbits in
$P^{\rm geom}_{\rm reg}$.  It is derived from the APF differential space and
its admissible-flow orbits, not assumed on the full reduced carrier.
\end{definition}

\subsection{What finite-ledger assumptions become geometrically}
The finite-ledger assumptions now divide cleanly.  Closed-world completeness
provides point separation and complete readout fibres.  Finite instruction
provides second countability.  Lossless postprocessing preserves those fibres,
and smooth saturation supplies the canonical Sikorski functional calculus.
Finite capacity provides finite-channel sufficiency; smooth local factorization
is the separate, explicitly named finite-differential-sufficiency gate.
Admissible continuation then supplies a family complete in the orbit-theorem sense of
flow-generating derivations, whose orbits are immersed smooth manifolds.
Constant rank supplies the quotient bundle used by the response geometry.

\subsection{Failure witnesses}
The differential-space formulation makes failure informative:
\begin{itemize}
\item a Stone-like zero-dimensional sector has records but no continuation
orbit of positive dimension;
\item a rank jump is an intrinsic singular stratum;
\item failure of finite-channel sufficiency signals saturation or unbounded
local capacity;
\item failure of finite differential sufficiency signals nonsmooth local
factorization even when the physical record remains complete;
\item a jump in the common derivation/cotangent dimension is a singularity
witness;
\item failure of concrete-model separation means only the numerical model, not
the universal physical response, is incomplete.
\end{itemize}

\section{Positive cost geometry and Lorentzian comparison geometry}
\label{sec:signature}

The pullback construction above is positive.  Physics also requires a signed causal comparison structure.  These are related structures, but they are not the same tensor obtained under different notation.

\begin{proposition}[Positive pullbacks cannot produce Lorentzian signature]
\label{prop:positive-pullback-no-lorentz}
Let $h$ be a positive semidefinite bilinear form on a process tangent space and $B:T_x\Mgeom\to T_{\mathcal N}\Pproc$ linear.  Then $B^*hB$ is positive semidefinite.  In particular, no construction of the form $T^*T$ with a positive ambient inner product can have Lorentzian index one.
\end{proposition}

\begin{proof}
For every $u$, $(B^*hB)(u,u)=h(Bu,Bu)\ge0$.  A Lorentzian form has vectors of both signs, so it cannot arise in this way.
\end{proof}

\begin{principle}[Causal cone from finite coordination]
If comparison continuations have finite maximum coordination throughput, tangent directions separate into causally admissible, null-boundary, and unreachable classes.  A smooth represented regime should encode these classes by a cone field $C_x\subset T_x\Mgeom$.
\end{principle}

\begin{principle}[Temporal orientation from irreversible record order]
Papers~1, 3, and 37 supply the distinction between coherent held identity and irreversible record-bearing order.  Only after a causal cone field has been independently constructed may Paper~9 represent that record preorder as a time orientation of the cone.  Record irreversibility therefore selects between the two components of an already available causal cone; it does not by itself derive Lorentzian signature or the cone.
\end{principle}

\begin{gate}[Signed Lorentz--Finsler comparison gate]
\label{gate:signed-comparison}
A signed comparison representation is available when there is a continuous two-homogeneous function
\[
        \Qcmp_x:T_x\Mgeom\longrightarrow\mathbb R
\]
that is $C^2$ on the relevant nonzero conic domains, has null set equal to the boundary of the finite-coordination cone, and whose vertical Hessian
\[
        g^{\rm cmp}_{x,v}
        :=\frac12D_v^2\Qcmp_x
\]
has Lorentzian index one.  If this Hessian is independent of $v$, then
\[
        \Qcmp_x(v)=q_x(v,v)
\]
for a Lorentzian bilinear form $q_x$.
\end{gate}

\begin{gate}[Cost--causal bridge]
The positive operational cost gauge $c_x^{\rm anc}$ and the signed comparison function $\Qcmp_x$ must be related by an explicit physical dictionary.  Possibilities include a common normalization on selected timelike or spacelike probes, a bound between implementation cost and $\sqrt{|\Qcmp|}$, or a variational relation.  No universal identity is assumed, especially on null directions.
\end{gate}

\begin{warning}[Why sectorwise polarization is insufficient]
Timelike and spacelike cones are not linear subspaces.  Imposing a parallelogram identity separately on positive- and negative-sign sectors does not by itself reconstruct the cross-sector bilinear pairing or a global Lorentzian form.  The signed two-homogeneous function or an equivalent cone Lagrangian must be defined on a common tangent domain before its Hessian can certify signature.
\end{warning}

\begin{gate}[Dimension gate]
The dimension gate closes when the number of independent local comparison modes is fixed by substrate capacity, representation stability, admissible interaction structure, and the comparison-response anchor.  A scalar-register lower bound may constrain one implementation architecture, but it is not by itself a theorem that the substrate tangent bundle has dimension four.  This draft continues to treat $(1+3)$ as an input to the weak-field application.
\end{gate}

\begin{audit}[Two geometric tensors]
In the quadratic branch the paper may carry both
\[
        \Gcost_x\quad\text{and}\quad q_x^{\rm cmp}.
\]
The first is positive and measures calibrated response effort.  The second is signed and encodes causal comparison.  They may share preferred directions or be linked by the recruitment dynamics, but identifying them requires a theorem beyond the pullback construction.
\end{audit}

\section{The geometric recruitment functional}

The geometric substrate is not a passive manifold.  It is the substrate that carries comparison continuations.  Matter and stress-energy distinctions recruit that substrate; the recruitment load changes the future cost of comparison.  The functional below is the geometric specialization of the recruitment cost functional of Papers~24/25.

\begin{definition}[Effective geometric recruitment load]
Let \(\chi_{\rm rec}(x)\) denote the effective accumulated recruitment load carried by the geometric substrate at \(x\).  The notation \(\chi\), rather than \(\rho\), is deliberate: the leading exterior point-mass profile is potential-like, scaling as \(1/r\), not as a volumetric density.
\end{definition}

\begin{definition}[Quadratic geometric recruitment functional]
In the linear-response regime around an unloaded substrate, a scalar recruitment load field \(\chi\) sourced by a mass-energy distinction \(\Jsrc\) is represented by
\[
        \Erec[\chi]
        =
        \frac12\int_{\Mgeom} \Kgeom^{ij}\,\partial_i\chi\,\partial_j\chi\,d^3x
        -
        \int_{\Mgeom} \Jsrc\,\chi\,d^3x
        + O(\chi^3),
\]
where \(\Kgeom^{ij}\) is the unloaded spatial comparison stiffness.  The cubic and higher terms encode nonlinear self-coupling and are not fixed in this draft.
\end{definition}

\begin{assumption}[Isotropic exterior kernel]
For a static, isolated, spherically symmetric source in the exterior linear-response regime, the unloaded spatial comparison stiffness is isotropic:
\[
        \Kgeom^{ij}=K_s\delta^{ij}.
\]
\end{assumption}

\begin{proposition}[Exterior point-source profile in the linear-response regime]
\label{prop:point-source-profile}
Under the quadratic recruitment functional and isotropic exterior kernel, a point source \(\Jsrc(x)=\mu M\delta^{(3)}(x)\) has an exterior equilibrium profile
\[
        \chi_{\rm rec}(r)=\frac{\mu M}{4\pi K_s r}
\]
up to the normalization convention for \(\mu\).
\end{proposition}

\begin{proof}
Varying \(\Erec\) gives the Euler-Lagrange equation
\[
        -\partial_i(\Kgeom^{ij}\partial_j\chi)=\Jsrc.
\]
With \(\Kgeom^{ij}=K_s\delta^{ij}\), this becomes
\[
        -K_s\nabla^2\chi=\mu M\delta^{(3)}(x).
\]
Using \(\nabla^2(1/r)=-4\pi\delta^{(3)}(x)\), the exterior solution is
\[
        \chi_{\rm rec}(r)=\frac{\mu M}{4\pi K_s r},
\]
with any additive constant fixed by the boundary condition \(\chi\to0\) as \(r\to\infty\).
\end{proof}

\begin{remark}[What this does and does not close]
\Cref{prop:point-source-profile} closes only the Green's-function step once the quadratic isotropic kernel is assumed.  It does not derive the kernel from APF primitives, does not derive the source normalization \(\mu\), and does not derive Newton's constant.  It is nevertheless useful: it shows that the required \(1/r\) exterior profile is the natural equilibrium profile of the simplest geometric recruitment functional.
\end{remark}

\section{Cost-curvature response}

\begin{principle}[Linear cost-curvature response]
In the exterior weak-load regime, the time-comparison cost-curvature responds linearly to the effective accumulated geometric recruitment load:
\[
        K_t(x)=K_{t0}\bigl(1-\alpha\chi_{\rm rec}(x)\bigr),
\]
where \(K_{t0}\) is the unloaded time-comparison stiffness and \(\alpha\) is the coupling between committed load and residual time-comparison stiffness.
\end{principle}

The sign is physical: substrate already committed to carrying recruitment load has reduced residual stiffness for additional identity-carrying time comparisons.  In represented geometry, this reduced stiffness appears as slower local proper-time accumulation relative to the unloaded coordinate normalization. In substrate-position-axis terms (Paper~0 \S\texttt{sec:substrate\_position\_axis}), the sign is the substrate's middle-regime saturation-direction physics at the geometric-substrate scale: the closer the substrate is to its saturation regime --- the more committed its capacity is to existing recruitment load --- the less residual middle-regime capacity remains for additional comparison continuations. The Bekenstein-bound saturation-regime cardinal datum (Paper~6 Supplement v2.10 \S sec:Lambda\_subspace, $S_{\rm dS} = 61\log 102$) is the global-scale saturation limit of the same response.

\begin{warning}[Status of the response law]
The response law is a structural commitment in this draft.  A closed derivation would require a nonlinear geometric recruitment functional, expansion around the loaded equilibrium, and proof that the leading correction to \(K_{t0}\) is exactly \(-\alpha\chi_{\rm rec}\).  That is one of the main open tasks of this paper.
\end{warning}

\subsection{Minimal perturbative response model}

The following lemma does not close the response law from APF primitives.  It formalizes the exact kind of load-probe coupling that would make the law unavoidable in a weak-load effective theory.

\begin{definition}[Time-comparison probe]
Let \(u\) be a small normalized probe field representing an infinitesimal rest-frame identity-comparison tick.  In the unloaded substrate its leading quadratic cost is
\[
        E_t[u]=\frac12\int K_{t0}\,u^2\,d\lambda,
\]
where \(d\lambda\) is the local comparison-volume element appropriate to the clock sector.
\end{definition}

\begin{assumption}[Bilinear load-probe coupling]
In the weak-load regime, the leading coupling between accumulated geometric recruitment load \(\chi\) and the time-comparison probe \(u\) is
\[
        E_t[u;\chi]
        =\frac12\int \bigl(K_{t0}-\lambda\chi(x)\bigr)u(x)^2\,d\lambda
        +O(u^3,\chi^2u^2),
\]
with \(\lambda>0\).  The minus sign states residual-stiffness depletion: already committed geometric load reduces available stiffness for an additional identity-comparison tick.
\end{assumption}

\begin{proposition}[Linear response from bilinear load-probe coupling]
\label{prop:response-from-coupling}
Under the bilinear load-probe coupling, the local time-comparison stiffness at fixed load is
\[
        K_t(x)=K_{t0}-\lambda\chi(x)+O(\chi^2),
\]
or equivalently
\[
        K_t(x)=K_{t0}\bigl(1-\alpha\chi(x)\bigr)+O(\chi^2),
        \qquad \alpha:=\lambda/K_{t0}.
\]
\end{proposition}

\begin{proof}
The local stiffness is the second variation of the probe cost with respect to the probe amplitude at the loaded equilibrium:
\[
        K_t(x)=\left.\frac{\delta^2 E_t[u;\chi]}{\delta u(x)^2}\right|_{u=0}.
\]
Taking the second variation of the displayed effective cost gives \(K_{t0}-\lambda\chi(x)\) to first order in \(\chi\).  Dividing by \(K_{t0}\) gives the normalized coefficient \(\alpha=\lambda/K_{t0}\).
\end{proof}

\begin{audit}[What the response lemma buys]
\Cref{prop:response-from-coupling} converts the response law from a naked ansatz into a precise perturbative target.  The remaining open work is now sharper: derive the negative bilinear load-probe coupling from the nonlinear geometric substrate functional, rather than merely assert the final linear response.
\end{audit}

\begin{definition}[Newtonian calibration]
In physical units, define the exterior point-mass calibration by
\[
        \alpha\chi_{\rm rec}(r)=\frac{2GM}{rc^2}.
\]
Equivalently, the combination \(\alpha\mu/(4\pi K_s)\) is fixed to \(2G/c^2\).  In this draft, \(G\) is a substrate-specific dimensional input, not an internally derived number.
\end{definition}

\begin{corollary}[Schwarzschild time factor as calibrated cost curvature]
\label{cor:schwarzschild-factor}
Under the point-source profile and Newtonian calibration,
\[
        K_t(r)=K_{t0}\left(1-\frac{2GM}{rc^2}\right)
\]
in the exterior linear-response regime.
\end{corollary}

\section{Proper time as normalized signed clock-sector comparison length}

\begin{definition}[Proper-time dictionary]
In the represented metric branch, the proper-time element for a static identity-preserving held continuation is the normalized signed clock-sector comparison length assigned by the represented causal kernel:
\[
        d\tau(r)=\sqrt{K_t(r)/K_{t0}}\,dt.
\]
This dictionary treats \(dt\) as the unloaded coordinate comparison parameter and \(d\tau\) as the local identity-continuation clock tick.
\end{definition}

\begin{remark}[Why the square root]
The local line element is quadratic in infinitesimal comparison displacement.  Cost-curvature multiplies the quadratic form; the represented cost length is therefore the square root of the curvature-normalized quadratic element.  This is a representation dictionary, not yet a derivation from a fully closed substrate action.
\end{remark}

\begin{warning}[Rest-frame bridge, not a signature derivation]
The proper-time dictionary is a calibration of the signed comparison kernel on static timelike probes.  It does not follow merely from the positive pullback response norm, and it does not identify general implementation cost with Lorentzian interval.  Null directions can have zero signed interval while carrying nonzero physical maintenance or realization cost.
\end{warning}

\begin{definition}[Clock sector]
The \emph{clock sector} of a represented geometry is the part of the comparison-cost kernel that controls static rest-frame identity continuations.  In a static coordinate representation it is encoded by a lapse factor
\[
        N(x):=\frac{d\tau(x)}{dt}.
\]
The redshift calculation constrains this sector only.
\end{definition}

\begin{definition}[Spatial comparison sector]
The \emph{spatial comparison sector} is the part of the comparison-cost kernel that controls endpoint comparison at fixed clock label.  In a weak-field isotropic representation one may write, schematically,
\[
        ds^2=-N(x)^2c^2dt^2 + S(x)^2\,d\mathbf x^2,
\]
where redshift fixes the lapse factor \(N\) but not the spatial scale factor \(S\).
\end{definition}

\begin{proposition}[Redshift does not fix full geometry]
\label{prop:redshift-not-full-geometry}
Any two static represented geometries with the same lapse factor \(N(r)=d\tau(r)/dt\) give the same static gravitational redshift between radii, even if their spatial comparison sectors differ.  Therefore a redshift recovery alone cannot establish full GR.
\end{proposition}

\begin{proof}
For static emitter and receiver, local frequency is inverse local proper period.  The same coordinate period is compared against the two local lapse factors, so
\[
        \frac{\omega_A}{\omega_B}=\frac{N(r_B)}{N(r_A)}.
\]
No spatial comparison factor enters this expression.  Hence all represented metrics with the same lapse produce the same redshift.
\end{proof}

\begin{proposition}[Leading-order redshift recovery from cost-curvature response]
\label{prop:redshift-recovery}
Consider a static, spherically symmetric recruitment load generated by a point mass \(M\), in the exterior linear-response regime.  Suppose
\[
        K_t(r)=K_{t0}\bigl(1-\alpha\chi_{\rm rec}(r)\bigr),
        \qquad
        \alpha\chi_{\rm rec}(r)=\frac{2GM}{rc^2},
\]
and suppose proper time is represented by normalized identity-continuation cost:
\[
        d\tau(r)=\sqrt{K_t(r)/K_{t0}}\,dt.
\]
Then
\[
        d\tau(r)=\sqrt{1-\frac{2GM}{rc^2}}\,dt.
\]
Consequently, for a photon emitted by a static observer at \(r_B\) and received by a static observer at \(r_A>r_B\),
\[
        \frac{\omega_A}{\omega_B}
        =
        \sqrt{\frac{1-2GM/r_Bc^2}{1-2GM/r_Ac^2}}.
\]
In the weak-field limit,
\[
        \frac{\omega_A}{\omega_B}
        \approx
        1-
        \frac{GM}{c^2}\left(\frac1{r_B}-\frac1{r_A}\right).
\]
\end{proposition}

\begin{proof}
The first displayed formula follows immediately from substituting the calibrated recruitment load into the cost-curvature response law and then applying the proper-time dictionary.  For a static emitter and receiver, the same coordinate-period label is compared against two different local proper-time rates.  Local frequency is inverse local proper period, hence
\[
        \omega_A=\omega_B\frac{d\tau(r_B)}{d\tau(r_A)}.
\]
Substitution gives the redshift ratio.  Expanding \(\sqrt{1-2GM/rc^2}\) to first order gives the weak-field expression.
\end{proof}

\begin{warning}[Boundary of the result]
\Cref{prop:redshift-recovery} is a conditional recovery, not a closed derivation of GR.  It assumes: (i) the linear cost-curvature response, (ii) the exterior point-source recruitment profile, (iii) the proper-time-as-normalized-cost dictionary, and (iv) \(G\) as the substrate coupling in physical units.  Deriving these from substrate primitives, extending the result beyond the exterior linear-response regime, and recovering the full Einstein equations are open.
\end{warning}

\section{No-smuggling audit for the redshift recovery}

The redshift recovery is useful only if its inputs are named without disguise.  The calculation has four layers.

\begin{longtable}{p{0.23\textwidth}p{0.16\textwidth}p{0.50\textwidth}}
\toprule
Layer & Status & Content \\
\midrule
Continuation grammar & \Status{Imported} & The judgment \(\G\entails_C D\goes E\) and the interpretation of geometry as comparison-continuation cost are imported from Paper~10. \\
Exterior profile & \Status{Cond./proved} & The \(1/r\) profile is proved here only after assuming an isotropic quadratic exterior kernel. \\
Clock response & \Status{Assumed} & \(K_t=K_{t0}(1-\alpha\chi)\) is a linear-response commitment. Its sign and linearity are physically motivated but not derived from a nonlinear substrate action. \\
Unit calibration & \Status{Input} & The equality \(\alpha\chi=2GM/rc^2\) imports the physical coupling \(G\). The calculation recovers the form of the redshift, not the numerical value of \(G\). \\
Proper-time dictionary & \Status{Representational} & \(d\tau=\sqrt{K_t/K_{t0}}dt\) is the metric-regime dictionary between normalized identity-continuation cost and clock time. \\
\bottomrule
\end{longtable}

\begin{audit}[What is actually earned]
The calculation earns the following claim: if a point mass produces a potential-like recruitment load, if time-comparison stiffness is reduced linearly by that load, and if proper time is normalized cost length, then the static clock-redshift sector has the Schwarzschild form at leading exterior order.  It does not earn the spatial metric, the equivalence principle, nonlinear field equations, or a derived value of \(G\).
\end{audit}

\begin{audit}[What would close the result]
To turn the recovery into a derivation, one must derive the isotropic Green's-function kernel, the universal response law, the proper-time dictionary, and the physical coupling normalization from APF substrate primitives.  These are not cosmetic gaps; they are the main mathematical work of Paper 9.
\end{audit}

\section{Newtonian acceleration as a cost-gradient effect}

The redshift calculation recovers the clock sector.  The next target is motion.  In represented weak-field geometry, slow test bodies accelerate along gradients of the same potential that controls \(g_{tt}\).  Cost-theoretically, this should emerge because least-cost identity-carrying worldlines bend toward lower continuation cost.

\begin{definition}[Effective cost potential]
Define the effective cost potential \(\Phi_C\) by
\[
        \frac{K_t(r)}{K_{t0}}=1+\frac{2\Phi_C(r)}{c^2}.
\]
For the calibrated point-source profile,
\[
        \Phi_C(r)=-\frac{GM}{r}.
\]
\end{definition}

\begin{proposition}[Newtonian acceleration under the geodesic-cost dictionary]
Assume the metric-representation gate closes and slow test-body worldlines extremize normalized identity-continuation cost.  In the weak-field, slow-velocity regime with
\[
        \frac{K_t}{K_{t0}}=1+\frac{2\Phi_C}{c^2},
\]
the represented coordinate acceleration is
\[
        \mathbf a=-\nabla\Phi_C.
\]
For \(\Phi_C=-GM/r\), this gives the inverse-square law.
\end{proposition}

\begin{proof}[Proof sketch]
Under the geodesic-cost dictionary, the action for a slow test body is proportional to
\[
        \int d\tau
        \approx
        \int \left(1+\frac{\Phi_C}{c^2}-\frac{v^2}{2c^2}\right)dt
\]
up to an irrelevant constant normalization and sign convention.  Extremizing is equivalent to extremizing the ordinary weak-field Lagrangian \(L=v^2/2-\Phi_C\), yielding \(\mathbf a=-\nabla\Phi_C\).  The step is conditional on the metric and geodesic-cost dictionaries.
\end{proof}

\begin{remark}[Why this matters]
The same scalar cost load controls both clock rate and slow-body acceleration.  This is the beginning of an equivalence-principle-like universality claim, but it is not yet a derivation of the equivalence principle.  To close that claim, the response law must be shown universal for all admissible rest-frame identity continuations, independent of matter composition.
\end{remark}


\section{Weak-field two-sector kernel and the missing spatial response}

The clock-sector recovery should be tested against the standard weak-field hierarchy.  The redshift calculation gives a lapse.  A gravitational theory also needs the spatial comparison response.  The next formal object is therefore not a scalar clock stiffness alone, but a two-sector normalized cost kernel.

\begin{definition}[Normalized two-sector weak-field cost kernel]
\label{def:two-sector-kernel}
In a static, isotropic, weak-field represented regime, write the normalized comparison-cost line element as
\[
        ds_C^2
        =-\Theta_t(r)c^2dt^2+\Theta_s(r)d\mathbf x^2+O(c^{-4}),
\]
where
\[
        \Theta_t(r):=\frac{K_t(r)}{K_{t0}},
        \qquad
        \Theta_s(r):=\frac{K_s(r)}{K_{s0}}.
\]
The redshift result fixes \(\Theta_t\).  It does not fix \(\Theta_s\).
\end{definition}

\begin{definition}[Weak-field cost potential]
For a static exterior point source, define the cost potential \(\Phi_C\) by
\[
        \Theta_t(r)=1+\frac{2\Phi_C(r)}{c^2}+O(c^{-4}).
\]
Under the Newtonian calibration of the scalar recovery,
\[
        \Phi_C(r)=-\frac{GM}{r}.
\]
\end{definition}

\begin{definition}[Spatial response parameter]
The weak-field spatial comparison response is defined by
\[
        \Theta_s(r)=1-2\gamma_C\frac{\Phi_C(r)}{c^2}+O(c^{-4}).
\]
The dimensionless coefficient \(\gamma_C\) is the APF/cost-theoretic counterpart of the standard PPN parameter \(\gamma_{\rm PPN}\).  General relativity has \(\gamma_C=1\).
\end{definition}

\begin{proposition}[Clock observables fix only the lapse]
\label{prop:clock-only-lapse}
Gravitational time dilation, static redshift, and Newtonian slow-body acceleration determine the leading clock potential \(\Phi_C\) but do not determine the spatial-response coefficient \(\gamma_C\).
\end{proposition}

\begin{proof}
The redshift and time-dilation formulae depend only on the lapse factor \(\Theta_t^{1/2}\).  The Newtonian slow-motion limit follows from the same lapse through the weak-field particle action.  The spatial factor \(\Theta_s\) enters null propagation and spatial-curvature observables at the same post-Newtonian order, but it does not enter the static redshift ratio and does not affect the lowest-order slow-velocity acceleration law.  Therefore the clock sector fixes \(\Phi_C\) and leaves \(\gamma_C\) free.
\end{proof}

\begin{gate}[Spatial-response / PPN-\(\gamma\) gate]
\label{gate:spatial-response}
To recover the full leading weak-field Schwarzschild sector, the geometric substrate must determine the spatial comparison response
\[
        \Theta_s(r)=1-2\frac{\Phi_C(r)}{c^2}+O(c^{-4}),
\]
that is,
\[
        \gamma_C=\gamma_{\rm PPN}=1.
\]
Until this gate closes, the paper has recovered the clock sector and Newtonian limit, but not light bending, Shapiro delay, or the full first-post-Newtonian exterior metric.

\end{gate}

\begin{gate}[Reciprocal clock-ruler response]
\label{gate:reciprocal-clock-ruler}
The spatial-response gate can be sharpened into a candidate closure principle.  In the exterior weak-field scalar-load regime, the time-comparison and spatial-comparison sectors are required to be reciprocal to first post-Newtonian order:
\[
        \Theta_s(r)\Theta_t(r)=1+O(c^{-4}),
\]
or equivalently
\[
        \Theta_s(r)=\Theta_t(r)^{-1}+O(c^{-4}).
\]
Physical reading: clock comparison and ruler comparison are conjugate faces of the same finite-coordination substrate.  A load that reduces the residual rate of identity-carrying clock continuation must increase the cost of spatial endpoint comparison by the inverse factor if the same local null coordination boundary is to be used by all admissible probes.  This is a gate, not a theorem: it must ultimately be derived from the microscopic comparison functional.
\end{gate}

\begin{proposition}[Reciprocal response implies the PPN spatial coefficient]
\label{prop:reciprocal-implies-gamma}
Assume the two-sector weak-field form
\[
        \Theta_t(r)=1+\frac{2\Phi_C(r)}{c^2}+O(c^{-4}),
        \qquad
        \Theta_s(r)=1-2\gamma_C\frac{\Phi_C(r)}{c^2}+O(c^{-4}).
\]
If the reciprocal clock-ruler gate holds, then
\[
        \gamma_C=1.
\]
\end{proposition}

\begin{proof}
Multiplying the two sector factors gives
\[
\begin{aligned}
        \Theta_s\Theta_t
        &=\left(1-2\gamma_C\frac{\Phi_C}{c^2}\right)
          \left(1+2\frac{\Phi_C}{c^2}\right)+O(c^{-4})\\
        &=1+2(1-\gamma_C)\frac{\Phi_C}{c^2}+O(c^{-4}).
\end{aligned}
\]
The reciprocal gate requires this product to be unity through order \(c^{-2}\).  For a nonzero exterior potential this forces \(1-\gamma_C=0\), hence \(\gamma_C=1\).
\end{proof}

\begin{audit}[Why naive same-sign scalar depletion fails]
A tempting first guess is that every comparison stiffness is depleted in the same scalar way:
\[
        \Theta_s(r)=\Theta_t(r)=1+\frac{2\Phi_C(r)}{c^2}+O(c^{-4}).
\]
Compared with the PPN form \(\Theta_s=1-2\gamma_C\Phi_C/c^2\), this gives \(\gamma_C=-1\).  Thus redshift would still be recovered from the clock sector, but the spatial-curvature contribution needed for the standard two-sector weak-field observables would have the wrong coefficient.  The geometric substrate cannot be modeled as uniform same-sign depletion of every comparison channel.  The spatial sector must be conjugate, tensorial, or otherwise constrained; this is precisely the content of \Cref{gate:reciprocal-clock-ruler} and \Cref{gate:spatial-response}.
\end{audit}

\subsection{Candidate derivation route: clock-ruler comparison-cell conservation}

The reciprocal clock-ruler condition should not remain a slogan.  The following definitions isolate a sharper operational certificate.  The certificate is still a gate, but it is weaker and more testable than simply assuming \(\gamma_C=1\): it says that, direction by direction, the elementary comparison cell built from one clock tick and one spatial endpoint ruler is preserved by a static scalar load to first post-Newtonian order.

\begin{definition}[Directionwise clock-ruler comparison cell]
Fix a represented spatial comparison direction \(e\) in the static weak-field regime.  The normalized clock-ruler comparison cell in that direction is
\[
        \mathcal A_e(r)
        :=
        \left(\frac{d\tau(r)}{dt}\right)
        \left(\frac{d\ell_e(r)}{dx_e}\right)
        =
        \sqrt{\Theta_t(r)\Theta_{s,e}(r)},
\]
where \(d\ell_e=\sqrt{\Theta_{s,e}}\,dx_e\) is the represented spatial comparison length in direction \(e\).  In the isotropic weak-field sector, \(\Theta_{s,e}=\Theta_s\) for all spatial directions.
\end{definition}

\begin{gate}[Directionwise comparison-cell conservation]
\label{gate:cell-conservation}
In an exterior, static, scalar-load regime with no shear, torsion, anisotropic substrate recruitment, or composition-dependent comparison channel, the elementary clock-ruler comparison cell is conserved to first post-Newtonian order:
\[
        \mathcal A_e(r)=1+O(c^{-4})
\]
for every represented spatial direction \(e\).

Physical reading: a scalar geometric load may redistribute comparison capacity between clock-rate continuation and endpoint-ruler continuation, but it does not create or destroy the local directionwise comparison cell.  Clock slowing and ruler stretching are conjugate responses of the same finite-coordination substrate.
\end{gate}

\begin{proposition}[Cell conservation implies reciprocal clock-ruler response]
\label{prop:cell-implies-reciprocal}
In the static isotropic two-sector weak-field kernel, directionwise comparison-cell conservation implies
\[
        \Theta_s(r)\Theta_t(r)=1+O(c^{-4}).
\]
Consequently, under the PPN expansion of \Cref{prop:reciprocal-implies-gamma}, it implies \(\gamma_C=1\).
\end{proposition}

\begin{proof}
In the isotropic sector, \(\mathcal A_e=\sqrt{\Theta_t\Theta_s}\) is independent of \(e\).  The cell-conservation gate gives \(\sqrt{\Theta_t\Theta_s}=1+O(c^{-4})\).  Squaring gives \(\Theta_t\Theta_s=1+O(c^{-4})\).  The implication \(\gamma_C=1\) then follows from \Cref{prop:reciprocal-implies-gamma}.
\end{proof}

\begin{audit}[What the cell-conservation route does not prove]
Cell conservation is not the local speed of light, coordinate invariance, or redshift in disguise.  Local null measurements give a universal measured speed once the metric representation is assumed; they do not by themselves fix \(\Theta_s\).  The new content is the claim that a static scalar recruitment load preserves the directionwise clock-ruler comparison cell at order \(c^{-2}\).  This can fail if the load creates anisotropic shear, if clock and ruler comparisons draw on independent substrate reservoirs, if matter species couple to different comparison kernels, or if the geometric load changes the local comparison-cell budget rather than redistributing it.  These are empirical and mathematical failure modes, not notational details.
\end{audit}


\subsection{Sharper candidate route: canonical comparison pairing}

The cell-conservation gate can itself be decomposed.  A comparison between events has two faces: a primal clock-ruler displacement and a dual frequency-wavevector probe.  The dimensionless action/phase pairing between these faces is the invariant bookkeeping object.  This subsection isolates the extra condition needed to make that pairing imply reciprocal clock-ruler response.

\begin{definition}[Local canonical comparison pair]
Fix a spatial comparison direction \(e\) in a represented local frame.  The primal comparison variables are
\[
        q_e=(\tau,\ell_e),
\]
where \(\tau\) is clock continuation and \(\ell_e\) is endpoint-ruler continuation in direction \(e\).  The dual probe variables are
\[
        p_e=(\omega,k_e),
\]
where \(\omega\) is frequency cost and \(k_e\) is directional wave-number cost.  The local comparison one-form is
\[
        \vartheta_e=-\omega\,d\tau+k_e\,d\ell_e .
\]
Equivalently, \(\vartheta_e/\hbar\) is the local dimensionless phase increment of a probe used to test the comparison.
\end{definition}

\begin{gate}[Canonical load transport]
\label{gate:canonical-load-transport}
A static scalar recruitment load acts canonically on the local comparison pair if, direction by direction, its transport map preserves the comparison one-form up to terms of order \(c^{-4}\):
\[
        T_\chi^*\vartheta_e=\vartheta_e+O(c^{-4}).
\]
For a diagonal weak-field response this means that if
\[
        d\tau\mapsto a_t\,d\tau,
        \qquad
        d\ell_e\mapsto a_s\,d\ell_e,
\]
then the dual variables transform inversely,
\[
        \omega\mapsto a_t^{-1}\omega,
        \qquad
        k_e\mapsto a_s^{-1}k_e,
\]
so that the phase/action pairing is not changed by relabeling the same comparison through the loaded substrate.
\end{gate}

\begin{audit}[Canonical pairing alone is not enough]
\label{audit:canonical-not-enough}
\Cref{gate:canonical-load-transport} preserves the action/phase pairing, but it does not by itself force \(a_ta_s=1\).  Any nonzero product \(a_ta_s\) can be paired with inverse dual scaling and still preserve \(\vartheta_e\).  Therefore the spatial coefficient \(\gamma_C=1\) requires one additional physical statement: the scalar load may exchange clock and ruler comparison capacity, but it may not dilate the elementary primal clock-ruler comparison cell at order \(c^{-2}\).
\end{audit}

\begin{gate}[Trace-free clock-ruler load transport]
\label{gate:trace-free-transport}
In the exterior static scalar-load regime, after removing coordinate relabelings and anisotropic shear, the primal action of the load on each local clock-ruler comparison plane is trace-free to first post-Newtonian order.  Equivalently, there exists a scalar load parameter \(\sigma(r)=O(c^{-2})\) such that
\[
        d\tau\mapsto e^{\sigma(r)}d\tau,
        \qquad
        d\ell_e\mapsto e^{-\sigma(r)}d\ell_e
\]
for every spatial direction \(e\), up to \(O(c^{-4})\).

Physical reading: a static scalar geometric load does not create an independent local comparison-cell resource.  It transfers comparison capacity between conjugate clock and ruler channels.  The trace part would represent creation or destruction of the directionwise comparison cell; the trace-free part represents redistribution.
\end{gate}

\begin{proposition}[Trace-free transport implies cell conservation]
\label{prop:trace-free-implies-cell}
Assume the represented two-sector weak-field kernel and trace-free clock-ruler load transport.  Then directionwise clock-ruler comparison-cell conservation holds:
\[
        \mathcal A_e(r)=1+O(c^{-4}).
\]
Consequently,
\[
        \Theta_s(r)\Theta_t(r)=1+O(c^{-4})
        \quad\text{and}\quad
        \gamma_C=1.
\]
\end{proposition}

\begin{proof}
By definition, the normalized clock factor is \(d\tau/dt=\sqrt{\Theta_t}\), and the normalized ruler factor in direction \(e\) is \(d\ell_e/dx_e=\sqrt{\Theta_{s,e}}\).  Trace-free transport gives these factors as \(e^{\sigma}\) and \(e^{-\sigma}\), up to \(O(c^{-4})\).  Hence
\[
        \mathcal A_e
        =\left(\frac{d\tau}{dt}\right)
         \left(\frac{d\ell_e}{dx_e}\right)
        =e^{\sigma}e^{-\sigma}+O(c^{-4})
        =1+O(c^{-4}).
\]
The remaining implications are exactly \Cref{prop:cell-implies-reciprocal,prop:reciprocal-implies-gamma}.
\end{proof}

\begin{corollary}[Nested spatial-gate chain]
\label{cor:nested-spatial-chain}
The leading spatial coefficient can be closed by any of the following increasingly primitive certificates:
\[
\begin{aligned}
        \text{trace-free clock-ruler load transport}
        &\Longrightarrow
        \text{directionwise cell conservation}\\
        &\Longrightarrow
        \Theta_s\Theta_t=1+O(c^{-4})\\
        &\Longrightarrow
        \gamma_C=1.
\end{aligned}
\]
Thus the next mathematical target is sharper than \(\gamma_C=1\): derive trace-free transport from the geometric recruitment functional, or exhibit a countermodel showing that APF permits a trace component.
\end{corollary}


\subsection{Variational route: fixed-cell response and the trace obstruction}

Trace-free transport should itself be generated by a local response principle, not merely imposed.  The cleanest way to state the target is to decompose the diagonal clock-ruler load response into its trace and trace-free parts.

\begin{definition}[Logarithmic clock-ruler response variables]
\label{def:log-response-variables}
Fix a represented spatial direction \(e\).  Write the local diagonal load response as
\[
        \frac{d\tau}{dt}=a_t=e^{\xi_t},
        \qquad
        \frac{d\ell_e}{dx_e}=a_s=e^{\xi_s}.
\]
Define the trace and exchange variables
\[
        u_e:=\xi_t+\xi_s,
        \qquad
        v_e:=\xi_t-\xi_s.
\]
Then \(u_e=\log\mathcal A_e\) is the logarithmic clock-ruler cell dilation, while \(v_e\) is the trace-free exchange between clock and ruler comparison channels.
\end{definition}

\begin{gate}[Fixed-cell variational response]
\label{gate:fixed-cell-variational-response}
In the exterior static scalar-load regime, after removing coordinate relabelings and shear, the local response of a directionwise clock-ruler comparison plane is governed to quadratic order by an effective potential of the form
\[
        \mathcal V_e(u_e,v_e;\chi)
        =
        \frac12 A_e u_e^2
        +\frac12 B_e\bigl(v_e-\beta_e\chi\bigr)^2
        +O\bigl(|u_e|^3+|v_e|^3+\chi^2|u_e|+\chi^2|v_e|\bigr),
\]
with \(A_e,B_e>0\).  The scalar load sources the exchange mode \(v_e\) but does not source the trace mode \(u_e\) at first order.

Physical reading: committed geometric recruitment may shift capacity between clock continuation and ruler comparison, but a scalar static load has no independent channel for creating or destroying the elementary clock-ruler comparison cell.  Any first-order source term proportional to \(\chi u_e\) would be a new local comparison-resource reservoir and must be separately justified.
\end{gate}

\begin{proposition}[Fixed-cell variational response implies trace-free transport]
\label{prop:fixed-cell-implies-trace-free}
Assume the fixed-cell variational response gate and \(\chi=O(c^{-2})\).  The local equilibrium response satisfies
\[
        u_e=O(\chi^2)=O(c^{-4}),
        \qquad
        v_e=\beta_e\chi+O(\chi^2).
\]
Consequently
\[
        \mathcal A_e=e^{u_e}=1+O(c^{-4}),
\]
so directionwise comparison-cell conservation holds.  In the isotropic two-sector weak-field regime this implies
\[
        \Theta_s\Theta_t=1+O(c^{-4}),
        \qquad
        \gamma_C=1.
\]
\end{proposition}

\begin{proof}
Equilibrium requires \(\partial\mathcal V_e/\partial u_e=0\) and \(\partial\mathcal V_e/\partial v_e=0\).  To first order,
\[
        \frac{\partial\mathcal V_e}{\partial u_e}=A_eu_e+O(\chi^2),
        \qquad
        \frac{\partial\mathcal V_e}{\partial v_e}=B_e(v_e-\beta_e\chi)+O(\chi^2).
\]
Since \(A_e,B_e>0\), the equilibrium solution is
\[
        u_e=O(\chi^2),
        \qquad
        v_e=\beta_e\chi+O(\chi^2).
\]
Because \(\chi=O(c^{-2})\), the trace is absent through order \(c^{-2}\).  Hence \(\mathcal A_e=e^{u_e}=1+O(c^{-4})\).  The final implications are \Cref{prop:cell-implies-reciprocal,prop:reciprocal-implies-gamma}.
\end{proof}

\subsection{Operator form of the no-trace target}
\label{sec:operator-no-trace}

The two-variable potential above is the scalar normal form of a more general loaded-equilibrium problem.  In the represented comparison-metric regime, let $q_0$ be the unloaded local comparison kernel and let $\mathsf S$ be the finite-dimensional gauge-fixed space of symmetric perturbations $h$ at that point.  Write
\[
        \operatorname{tr}_{q_0}h:=q_0^{ab}h_{ab},
        \qquad
        \Ptr h:=\frac{1}{m}(\operatorname{tr}_{q_0}h)q_0,
        \qquad
        \Ptf:=I-\Ptr,
\]
where $m=\dim\Mgeom$ in the represented regime.

\begin{assumption}[Gauge-fixed quadratic recruitment response]
\label{ass:gauge-fixed-response-hessian}
Around $q_0$, the local recruitment functional has expansion
\[
\begin{aligned}
\mathcal E_{\rm rec}[q_0+h,\chi]
={}&\mathcal E_0
+\frac12\langle h,\mathcal K h\rangle
-\chi\langle\mathcal J,h\rangle \\
&+O\!\left(\|h\|^3+\chi\|h\|^2+\chi^2\|h\|\right),
\end{aligned}
\]
where $\mathcal K:\mathsf S\to\mathsf S$ is self-adjoint and invertible after coordinate-gauge zero modes are removed, and $\mathcal J\in\mathsf S$ is the scalar-load source.
\end{assumption}

\begin{theorem}[Linear response and the exact no-trace operator]
\label{thm:no-trace-operator}
Under \Cref{ass:gauge-fixed-response-hessian}, the equilibrium perturbation is
\[
        h_\chi=\chi\,\mathcal K^{-1}\mathcal J+O(\chi^2).
\]
Its first-order comparison-volume variation is
\[
        \delta_\chi\log\sqrt{|\det(q_0+h_\chi)|}
        =\frac{\chi}{2}\operatorname{tr}_{q_0}
        (\mathcal K^{-1}\mathcal J)+O(\chi^2).
\]
Consequently the no-trace gate is exactly
\[
        \boxed{\Ptr\mathcal K^{-1}\mathcal J=0}.
\]
\end{theorem}

\begin{proof}
The stationary equation is
\[
        \mathcal K h_\chi-\chi\mathcal J
        +O(\|h_\chi\|^2+\chi\|h_\chi\|+\chi^2)=0.
\]
Invertibility and the implicit-function expansion give $h_\chi=\chi\mathcal K^{-1}\mathcal J+O(\chi^2)$.  The determinant variation formula gives
\[
        \log\sqrt{|\det(q_0+h)|}
        =\log\sqrt{|\det q_0|}
        +\frac12\operatorname{tr}_{q_0}h+O(\|h\|^2).
\]
Substitution proves the volume formula.  Vanishing of the trace component is equivalent to $\Ptr\mathcal K^{-1}\mathcal J=0$.
\end{proof}

\begin{corollary}[Source-selection route]
\label{cor:source-selection-no-trace}
If
\[
        \Ptr\mathcal J=0
\]
and $\mathcal K$ preserves the trace-free subspace, equivalently
\[
        \Ptr\mathcal K\Ptf=0=\Ptf\mathcal K\Ptr,
\]
then $\Ptr\mathcal K^{-1}\mathcal J=0$ and the linear response is trace-free.
\end{corollary}

\begin{corollary}[Constrained fixed-volume route]
\label{cor:constrained-no-trace}
If the admissible variation space is constrained by
\[
        \operatorname{tr}_{q_0}h=0
\]
through a Lagrange multiplier or a primitive fixed-comparison-capacity condition, then the equilibrium response is trace-free by construction.  The remaining physical obligation is to derive the constraint from admissibility rather than impose it as a coordinate convention.
\end{corollary}

\begin{audit}[Relation to the scalar potential]
In a clock--ruler plane, $u_e$ is the trace coordinate and $v_e$ the trace-free exchange coordinate.  The condition that no term $\chi u_e$ appear is the two-dimensional source-selection statement $\Ptr\mathcal J=0$.  The stronger invariant target is $\Ptr\mathcal K^{-1}\mathcal J=0$, because a trace-free source can still generate trace if the Hessian mixes the trace and trace-free sectors.
\end{audit}

\begin{definition}[Trace-obstruction coefficient]
\label{def:trace-obstruction-coefficient}
If the local response potential contains a first-order trace source,
\[
        \mathcal V_e(u_e,v_e;\chi)
        =
        \frac12 A_eu_e^2+\frac12 B_e(v_e-\beta_e\chi)^2
        -\lambda_e\chi u_e+O(3),
\]
then the equilibrium trace is
\[
        u_e=\eta_e\chi+O(\chi^2),
        \qquad
        \eta_e:=\lambda_e/A_e.
\]
The coefficient \(\eta_e\) is the first-order trace obstruction.
\end{definition}

\begin{proposition}[Trace obstruction is the PPN-\(\gamma\) obstruction]
\label{prop:trace-obstruction-gamma}
In the isotropic two-sector weak-field regime, suppose the logarithmic cell dilation has the expansion
\[
        u(r)=\zeta_C\frac{\Phi_C(r)}{c^2}+O(c^{-4}).
\]
Then
\[
        \gamma_C=1-\zeta_C.
\]
Thus the GR spatial coefficient \(\gamma_C=1\) is equivalent to vanishing first-order trace dilation \(\zeta_C=0\).
\end{proposition}

\begin{proof}
By definition,
\[
        e^{2u}=\Theta_t\Theta_s.
\]
Using the two-sector weak-field expansion,
\[
\begin{aligned}
        \Theta_t\Theta_s
        &=\left(1+2\frac{\Phi_C}{c^2}\right)
          \left(1-2\gamma_C\frac{\Phi_C}{c^2}\right)+O(c^{-4})\\
        &=1+2(1-\gamma_C)\frac{\Phi_C}{c^2}+O(c^{-4}).
\end{aligned}
\]
Meanwhile \(e^{2u}=1+2\zeta_C\Phi_C/c^2+O(c^{-4})\).  Equating coefficients gives \(\zeta_C=1-\gamma_C\), hence \(\gamma_C=1-\zeta_C\).
\end{proof}

\begin{audit}[The remaining obstruction]
The spatial gate has now been reduced to a concrete question: does static scalar geometric recruitment source the trace mode \(u_e\) of the local clock-ruler response?  If not, fixed-cell variational response gives trace-free transport and closes \(\gamma_C=1\).  If yes, the sourced trace coefficient \(\zeta_C\) shifts the PPN spatial coefficient by \(\gamma_C=1-\zeta_C\).  The obstruction is no longer vague.  It is the existence of a first-order \(\chi u_e\) term in the effective local response potential.
\end{audit}


\subsection{Generator form: the obstruction is a \texorpdfstring{$GL(2)$}{GL(2)} trace}

The variational statement above is useful, but the same obstruction can be stated more invariantly.  A weak static load acts locally on each clock-ruler comparison plane by a linear response map.  The determinant of this map is the comparison-cell factor.  The trace of its infinitesimal generator is therefore exactly the first-order cell-dilation obstruction.

\begin{definition}[Clock-ruler response map]
\label{def:clock-ruler-response-map}
Fix a represented spatial direction \(e\).  After removing coordinate relabelings and shear, write the loaded clock-ruler increments as
\[
        \begin{pmatrix} d\tau \\ d\ell_e \end{pmatrix}
        =R_e(\hat\chi)
        \begin{pmatrix} d\tau_0 \\ d\ell_{e0} \end{pmatrix},
        \qquad
        \hat\chi(r):=\frac{\Phi_C(r)}{c^2}=O(c^{-2}),
\]
where \(R_e(0)=I\) and \(R_e(\hat\chi)\in GL^+(2,\mathbb R)\) in the represented smooth regime.  Its infinitesimal generator is
\[
        L_e:=\left.\frac{dR_e}{d\hat\chi}\right|_{0},
        \qquad
        R_e(\hat\chi)=I+\hat\chi L_e+O(\hat\chi^2).
\]
The trace component
\[
        \theta_e:=\operatorname{tr}L_e
\]
is the infinitesimal comparison-cell dilation rate.  The traceless part \(L_e-\frac12\theta_e I\) is the exchange/shear part of the clock-ruler response.
\end{definition}

\begin{proposition}[Generator trace equals trace obstruction]
\label{prop:generator-trace-obstruction}
In the diagonal isotropic weak-field sector, the logarithmic cell dilation satisfies
\[
        u_e(r)=\theta_e\frac{\Phi_C(r)}{c^2}+O(c^{-4}).
\]
Consequently
\[
        \zeta_C=\theta_e,
        \qquad
        \gamma_C=1-\theta_e.
\]
Thus the GR spatial coefficient \(\gamma_C=1\) is equivalent, at this order, to the infinitesimal load generator being traceless:
\[
        \operatorname{tr}L_e=0.
\]
\end{proposition}

\begin{proof}
The comparison-cell factor is
\[
        \mathcal A_e=\det R_e(\hat\chi)
        =1+\hat\chi\operatorname{tr}L_e+O(\hat\chi^2).
\]
Since \(u_e=\log\mathcal A_e\), we obtain
\[
        u_e=\hat\chi\operatorname{tr}L_e+O(\hat\chi^2)
        =\theta_e\Phi_C/c^2+O(c^{-4}).
\]
Comparison with \Cref{prop:trace-obstruction-gamma} gives \(\zeta_C=\theta_e\) and hence \(\gamma_C=1-\theta_e\).
\end{proof}


\subsection{Trace as comparison-volume variation}
\label{sec:trace-volume}

The $GL(2)$ trace has an exact geometric interpretation once a represented signed comparison kernel is available on the local clock--ruler plane.

\begin{definition}[Loaded clock--ruler comparison kernel]
Let $q_{e,0}$ be the unloaded nondegenerate comparison kernel on the clock--ruler plane.  Represent the loaded kernel on the unloaded increment space by
\[
        q_e(\hat\chi):=R_e(\hat\chi)^*q_{e,0}
        =R_e(\hat\chi)^{\mathsf T}q_{e,0}R_e(\hat\chi).
\]
\end{definition}

\begin{theorem}[Trace--volume identity]
\label{thm:trace-volume}
For $R_e(\hat\chi)=I+\hat\chi L_e+O(\hat\chi^2)$,
\[
        \boxed{
        \left.\frac{d}{d\hat\chi}
        \log\sqrt{|\det q_e(\hat\chi)|}
        \right|_{\hat\chi=0}
        =\operatorname{tr}L_e
        }.
\]
Thus the infinitesimal generator trace is precisely the first variation of the local comparison-volume element.
\end{theorem}

\begin{proof}
The determinant transforms as
\[
        \det q_e(\hat\chi)
        =\det(R_e)^2\det q_{e,0}.
\]
Because $R_e(0)=I$ and $R_e\in GL^+(2,\mathbb R)$ locally,
\[
        \log\sqrt{|\det q_e|}
        =\log\det R_e+\log\sqrt{|\det q_{e,0}|}.
\]
Differentiating at zero and using $\frac{d}{d\hat\chi}\log\det R_e|_0=\operatorname{tr}L_e$ proves the identity.
\end{proof}

\begin{corollary}[The PPN obstruction is the scalar volume mode]
\label{cor:ppn-volume-mode}
After clock calibration,
\[
        \left.\frac{d}{d\hat\chi}
        \log\sqrt{|\det q_e(\hat\chi)|}
        \right|_{0}
        =1-\gamma_C.
\]
Hence $\gamma_C=1$ if and only if the static scalar load has no first-order local comparison-volume mode.
\end{corollary}

\subsection{Gauge-reduced scalar normal form}

The generator formulation still permits too many apparent degrees of freedom.  Some are physical, but some are only coordinate choice, anisotropic shear, or time-space mixing inappropriate to the static scalar exterior sector.  The next reduction is to put the local response into normal form.

\begin{definition}[Gauge-reduced scalar response sector]
\label{def:gauge-reduced-scalar-sector}
Fix a static, spherically symmetric exterior weak-field load.  A local clock-ruler response map is in the gauge-reduced scalar sector when, after quotienting coordinate relabelings, anisotropic shear, and parity-forbidden time-space mixing, its infinitesimal generator is diagonal and direction-independent:
\[
        L_{\rm sc}
        =
        \begin{pmatrix}
        \lambda_t & 0\\
        0 & \lambda_s
        \end{pmatrix}
        =\frac{\theta}{2}I+\sigma H,
        \qquad
        H:=
        \begin{pmatrix}
        1&0\\
        0&-1
        \end{pmatrix}.
\]
Here
\[
        \theta:=\lambda_t+\lambda_s=\operatorname{tr}L_{\rm sc},
        \qquad
        \sigma:=\frac{\lambda_t-\lambda_s}{2}.
\]
The trace \(\theta\) is the comparison-cell dilation mode; the exchange coefficient \(\sigma\) transfers comparison scale between clock and ruler channels without changing the cell determinant.
\end{definition}

\begin{proposition}[Clock calibration leaves only the trace obstruction]
\label{prop:normal-form-trace-obstruction}
In the gauge-reduced scalar sector, suppose the clock response is calibrated by the redshift recovery,
\[
        \sqrt{\Theta_t}=1+\frac{\Phi_C}{c^2}+O(c^{-4}),
\]
and write the spatial response as
\[
        \sqrt{\Theta_s}=1-\gamma_C\frac{\Phi_C}{c^2}+O(c^{-4}).
\]
Then the normal-form generator has
\[
        \lambda_t=1,
        \qquad
        \lambda_s=-\gamma_C,
\]
so that
\[
        \theta=1-
        \gamma_C,
        \qquad
        \sigma=\frac{1+
        \gamma_C}{2}.
\]
Therefore, once the clock sector is calibrated, the entire first-order spatial-response problem is the trace problem:
\[
        \gamma_C=1
        \quad\Longleftrightarrow\quad
        \theta=0.
\]
\end{proposition}

\begin{proof}
By \Cref{def:clock-ruler-response-map}, the loaded increments are
\[
        d\tau=(1+\hat\chi\lambda_t)d\tau_0+O(\hat\chi^2),
        \qquad
        d\ell=(1+\hat\chi\lambda_s)d\ell_0+O(\hat\chi^2),
\]
with \(\hat\chi=\Phi_C/c^2\).  Matching these expressions to the displayed clock and ruler factors gives \(\lambda_t=1\) and \(\lambda_s=-\gamma_C\).  The formulae for \(\theta\) and \(\sigma\) follow from \Cref{def:gauge-reduced-scalar-sector}.  Hence \(\gamma_C=1\) iff \(\theta=0\).
\end{proof}

\begin{corollary}[Same-sign depletion failure in normal form]
\label{cor:same-sign-failure-normal-form}
A same-sign scalar depletion of clock and ruler increments has \(\lambda_t=\lambda_s=1\).  With the clock sector calibrated, this implies
\[
        \gamma_C=-1,
\]
so it reproduces the redshift sign but fails the spatial weak-field sector.
\end{corollary}

\begin{proof}
If \(\lambda_t=1\) and \(\lambda_s=1\), then \(-\gamma_C=1\), hence \(\gamma_C=-1\).  Equivalently, the trace is \(\theta=2\), so \(\gamma_C=1-\theta=-1\).
\end{proof}

\begin{proposition}[Unique scalar closure compatible with calibrated weak-field GR]
\label{prop:unique-scalar-closure}
Within the gauge-reduced scalar sector with the clock coefficient fixed by redshift, the following are equivalent to first post-Newtonian order:
\[
\begin{aligned}
        L_{\rm sc}\in\mathfrak{sl}(2,\mathbb R)
        &\Longleftrightarrow
        L_{\rm sc}=H \\
        &\Longleftrightarrow
        R_{\rm sc}\in SL(2,\mathbb R)+O(c^{-4}) \\
        &\Longleftrightarrow
        \Theta_s=\Theta_t^{-1}+O(c^{-4}) \\
        &\Longleftrightarrow
        \gamma_C=1.
\end{aligned}
\]
Thus the scalar weak-field closure is unique: the load must act as a trace-free clock-ruler exchange, not as a common dilation.
\end{proposition}

\begin{proof}
Clock calibration gives \(\lambda_t=1\).  Tracelessness gives \(\lambda_s=-1\), hence \(L_{\rm sc}=H\).  This is equivalent to \(\det R_{\rm sc}=1+O(c^{-4})\), since the determinant changes to first order by \(\operatorname{tr}L_{\rm sc}\hat\chi\).  The determinant condition is \(\Theta_s\Theta_t=1+O(c^{-4})\), which is the same as \(\Theta_s=\Theta_t^{-1}+O(c^{-4})\).  By \Cref{prop:normal-form-trace-obstruction}, this is equivalent to \(\gamma_C=1\).
\end{proof}

\begin{theorem}[Scalar weak-field closure package]
\label{thm:scalar-closure-package}
Assume the static isotropic exterior weak-field sector, the metric-representation gate, universal matter coupling to the same comparison kernel, and the clock calibration supplied by the redshift recovery.  After coordinate gauge, anisotropic shear, and parity-forbidden time-space mixing are removed, the following statements are equivalent to first post-Newtonian order:
\[
\begin{aligned}
        \text{no scalar comparison-cell dilation}
        &\Longleftrightarrow
        \operatorname{tr}L_{\rm sc}=0 \\
        &\Longleftrightarrow
        L_{\rm sc}\in\mathfrak{sl}(2,\mathbb R) \\
        &\Longleftrightarrow
        \det R_{\rm sc}=1+O(c^{-4}) \\
        &\Longleftrightarrow
        \left.\dfrac{d}{d\hat\chi}
        \log\sqrt{|\det q_{\rm sc}(\hat\chi)|}\right|_{0}=0 \\
        &\Longleftrightarrow
        \Theta_s\Theta_t=1+O(c^{-4}) \\
        &\Longleftrightarrow
        \gamma_C=1.
\end{aligned}
\]
If these equivalent conditions hold, the conditional weak-field exterior recovery theorem applies.  If they fail with trace \(\theta\), then
\[
        \gamma_C=1-\theta,
\]
so two-sector weak-field observations measure the scalar comparison-cell dilation directly.
\end{theorem}

\begin{proof}
Clock calibration gives \(\lambda_t=1\).  In the gauge-reduced scalar sector, \Cref{prop:normal-form-trace-obstruction} gives \(\theta=1-\gamma_C\), so \(\gamma_C=1\) iff \(\theta=0\).  The equivalence between vanishing trace, membership in \(\mathfrak{sl}(2,\mathbb R)\), first-order unimodularity of \(R_{\rm sc}\), and preservation of the normalized clock-ruler cell is the determinant identity for a first-order \(GL(2)\) response.  In the isotropic two-sector kernel this determinant condition is \(\Theta_s\Theta_t=1+O(c^{-4})\).  The conditional weak-field recovery then follows from \Cref{thm:conditional-weak-field}.  If \(\theta\neq0\), \Cref{prop:trace-obstruction-gamma} gives the shifted coefficient \(\gamma_C=1-\theta\).
\end{proof}

\begin{audit}[What the scalar package reduces]
The spatial problem is no longer an unconstrained metric-fitting problem.  In the static scalar exterior sector, after gauge and shear are removed, redshift fixes \(\lambda_t\).  The only remaining first-order scalar ambiguity is \(\theta=\operatorname{tr}L_{\rm sc}\).  Proving the APF substrate has no independent scalar comparison-cell dilation is therefore enough to close the leading two-sector weak-field geometry.
\end{audit}

\begin{gate}[Unimodular scalar load response]
\label{gate:unimodular-load-response}
Static scalar geometric recruitment acts on each local clock-ruler comparison plane by a unimodular response after coordinate gauge, anisotropic shear, and composition-dependent channels are removed:
\[
        \det R_e(\hat\chi)=1+O(c^{-4}),
        \qquad\text{equivalently}\qquad
        L_e\in\mathfrak{sl}(2,\mathbb R).
\]
Physical reading: scalar recruitment may redistribute comparison capacity between clock and ruler channels, but it does not create or destroy the elementary local comparison cell.  A failure of this gate is not a small technical issue; it is an additional scalar trace mode in the weak-field geometry.
\end{gate}

\begin{proposition}[Unimodular response closes the spatial coefficient]
\label{prop:unimodular-implies-gamma}
Assume the two-sector weak-field kernel and the unimodular scalar load response gate.  Then
\[
        \Theta_s\Theta_t=1+O(c^{-4}),
        \qquad
        \gamma_C=1.
\]
\end{proposition}

\begin{proof}
The determinant of the clock-ruler response map is the normalized comparison-cell factor:
\[
        \det R_e=\mathcal A_e
        =\left(\frac{d\tau}{dt}\right)
         \left(\frac{d\ell_e}{dx_e}\right)
        =\sqrt{\Theta_t\Theta_{s,e}}.
\]
Unimodularity gives \(\mathcal A_e=1+O(c^{-4})\), hence \(\Theta_t\Theta_{s,e}=1+O(c^{-4})\).  In the isotropic sector this is \(\Theta_t\Theta_s=1+O(c^{-4})\), and \Cref{prop:reciprocal-implies-gamma} gives \(\gamma_C=1\).
\end{proof}

\begin{corollary}[Equivalence of the spatial closure targets]
\label{cor:spatial-target-equivalence}
In the isotropic weak-field regime, the following are equivalent to first post-Newtonian order:
\[
\begin{aligned}
        \operatorname{tr}L_e=0
        &\Longleftrightarrow
        \det R_e=1+O(c^{-4}) \\
        &\Longleftrightarrow
        \mathcal A_e=1+O(c^{-4}) \\
        &\Longleftrightarrow
        \Theta_s\Theta_t=1+O(c^{-4}) \\
        &\Longleftrightarrow
        \gamma_C=1.
\end{aligned}
\]
The spatial gate can therefore be pursued as any one of four equivalent problems: rule out the trace generator, prove unimodular response, prove comparison-cell conservation, or derive reciprocal clock-ruler response.
\end{corollary}

\begin{audit}[Operational diagnostic for the trace mode]
If weak-field tests determine an effective PPN coefficient \(\gamma_{\rm obs}\), the APF trace mode is read off as
\[
        \theta_e=\zeta_C=1-\gamma_{\rm obs}.
\]
Thus the observational role of light bending, Shapiro delay, and other two-sector probes is not merely to test GR.  In the cost language, they bound or detect the first-order determinant of the local clock-ruler response map.  Redshift alone cannot do this, because it probes only the clock projection.
\end{audit}

\begin{proposition}[Nonzero trace is a scalar comparison mode]
\label{prop:nonzero-trace-scalar-mode}
Within the two-sector weak-field ansatz, a nonzero first-order trace generator \(\theta_e\neq0\) is observationally equivalent, at this order, to adding an independent scalar comparison mode that changes the spatial PPN coefficient while leaving the redshift sector fixed.
\end{proposition}

\begin{proof}
The redshift sector fixes \(\Theta_t=1+2\Phi_C/c^2+O(c^{-4})\).  A nonzero trace generator changes the product \(\Theta_t\Theta_s\) by
\[
        \Theta_t\Theta_s
        =1+2\theta_e\frac{\Phi_C}{c^2}+O(c^{-4}),
\]
which by \Cref{prop:generator-trace-obstruction} shifts \(\gamma_C\) to \(1-\theta_e\).  Since \(\Theta_t\) is unchanged, the new degree of freedom is invisible to clock-sector redshift but visible to two-sector spatial observables.  That is exactly the weak-field role of an additional scalar comparison mode.
\end{proof}

\begin{remark}[Why this is a strengthening]
The variational obstruction can be expressed invariantly as a generator question: is scalar load transport \(GL(2)\) with nonzero trace, or \(SL(2)\) to first order?  This is the cleanest current statement of the spatial-response problem.
\end{remark}

\begin{corollary}[A sharper spatial-response target]
\label{cor:sharper-target}
The weak-field spatial sector can be closed by proving the unimodular scalar load response gate \Cref{gate:unimodular-load-response}, or equivalently any of the targets listed in \Cref{cor:spatial-target-equivalence}.  One convenient proof chain is
\[
\begin{aligned}
        \operatorname{tr}L_e=0
        &\Longrightarrow
        \det R_e=1+O(c^{-4}) \\
        &\Longrightarrow
        \Theta_s\Theta_t=1+O(c^{-4}) \\
        &\Longrightarrow
        \gamma_C=1 \\
        &\Longrightarrow
        \text{leading two-sector weak-field GR recovery}.
\end{aligned}
\]
Thus the next target is no longer the bare coefficient \(\gamma_C\), but the underlying no-trace / unimodular comparison-load certificate.
\end{corollary}

\begin{assumption}[Two-sector weak-field response]
\label{ass:two-sector-response}
In the exterior weak-field regime generated by a static point source, the represented comparison kernel has
\[
        \Theta_t(r)=1+\frac{2\Phi_C(r)}{c^2}+O(c^{-4}),
        \qquad
        \Theta_s(r)=1-2\gamma_C\frac{\Phi_C(r)}{c^2}+O(c^{-4}).
\]
The clock calculation gives the first relation conditionally.  The second relation is the spatial-response target.
\end{assumption}

\begin{theorem}[Conditional weak-field exterior recovery]
\label{thm:conditional-weak-field}
Assume the metric-representation gate, the exterior point-source profile, the Newtonian calibration \(\Phi_C=-GM/r\), the proper-time dictionary, universal matter coupling to the same comparison kernel, and the spatial-response gate \(\gamma_C=1\).  Then the represented exterior weak-field cost kernel is
\[
        ds_C^2
        =-\left(1-\frac{2GM}{rc^2}\right)c^2dt^2
        +\left(1+\frac{2GM}{rc^2}\right)d\mathbf x^2
        +O(c^{-4}),
\]
which is the first-post-Newtonian Schwarzschild exterior metric in isotropic weak-field form.  Consequently the framework conditionally recovers, at leading weak-field order, gravitational redshift, clock-height tests, Newtonian acceleration, light bending, and Shapiro time delay.
\end{theorem}

\begin{proof}
The clock-sector recovery and calibration give
\[
        \Theta_t(r)=1-\frac{2GM}{rc^2}+O(c^{-4}).
\]
The spatial-response gate with \(\gamma_C=1\) gives
\[
        \Theta_s(r)=1+\frac{2GM}{rc^2}+O(c^{-4}).
\]
Substitution into \Cref{def:two-sector-kernel} gives the displayed metric kernel.  In standard post-Newtonian bookkeeping, this pair of coefficients is exactly the weak-field Schwarzschild exterior through the leading order controlling the listed observables.  The recovery is conditional because the metric-representation gate, universal matter coupling, and either the spatial-response gate or its sharper reciprocal clock-ruler route are not yet derived from substrate primitives.
\end{proof}


\begin{corollary}[No-trace scalar response is sufficient for the leading weak-field package]
\label{cor:no-trace-suffices-weak-field}
Under the hypotheses of \Cref{thm:conditional-weak-field}, the explicit spatial-response assumption \(\gamma_C=1\) may be replaced by any one of the equivalent scalar-closure conditions in \Cref{thm:scalar-closure-package}.  In particular, if the gauge-reduced scalar load generator is traceless,
\[
        \operatorname{tr}L_{\rm sc}=0,
\]
then the conditional first-post-Newtonian exterior recovery follows.
\end{corollary}

\begin{proof}
By \Cref{thm:scalar-closure-package}, vanishing trace is equivalent to \(\gamma_C=1\) in the static isotropic scalar exterior sector after clock calibration.  Substituting this equivalent gate for the spatial-response assumption gives exactly the hypotheses of \Cref{thm:conditional-weak-field}.
\end{proof}

\begin{corollary}[Observable dependency split]
\label{cor:observable-dependency-split}
Within the two-sector weak-field kernel, redshift and clock-height tests are clock-sector observables, while light bending and Shapiro delay are two-sector observables.  Therefore any APF claim to recover leading weak-field GR must close both the clock-response gate and the spatial-response gate.
\end{corollary}

\begin{audit}[No-smuggling table for the weak-field suite]
\begin{longtable}{p{0.22\textwidth}p{0.16\textwidth}p{0.16\textwidth}p{0.33\textwidth}}
\toprule
Observable & Needs \(\Theta_t\) & Needs \(\Theta_s\) & Present status \\
\midrule
Static redshift & Yes & No & Recovered conditionally from clock response. \\
Clock-height tests & Yes & No & Same gate as redshift, with local calibration. \\
Newtonian acceleration & Yes & No at leading slow-body order & Conditional on geodesic-cost dictionary and universal matter coupling. \\
Light bending & Yes & Yes & Open until \(\gamma_C=1\) is derived. \\
Shapiro delay & Yes & Yes & Open until null comparison uses the two-sector kernel with \(\gamma_C=1\). \\
Perihelion precession & Yes & Yes plus next-order terms & Partly first-PN; full recovery needs nonlinear corrections. \\
Frame dragging & Needs \(h_{0i}\) & Yes & Deferred to rotating-source tensor sector. \\
Gravitational radiation & Tensor dynamics & Tensor dynamics & Deferred to dynamical self-coupled geometric substrate. \\
\bottomrule
\end{longtable}
\end{audit}

\section{From response cost geometry to tensorial comparison response}
\label{sec:tensorial-response}

The scalar calculation used one effective load $\chi$ and one clock stiffness.  The intrinsic construction now shows that a tensorial theory may carry two different quadratic objects:
\[
        \Gcost_{ab}
        \qquad\text{and}\qquad
        q^{\rm cmp}_{ab}.
\]
The first, when it exists, is the positive metric induced by calibrated response effort.  The second is the signed comparison kernel supplied by the causal/signature gate.  Their relation is a bridge theorem, not an identity inherited from notation.

\begin{definition}[Tensorial signed comparison kernel]
A tensorial signed comparison kernel is a smooth nondegenerate field $q^{\rm cmp}_{ab}$ whose quadratic function
\[
        \Qcmp_x(dx)=q^{\rm cmp}_{ab}(x)\,dx^a dx^b
\]
represents the causal comparison structure.  When the signature gate closes it has Lorentzian signature.  The positive implementation effort of the same displacement is measured separately by $\Gcost$ or by the underlying Finsler gauge.
\end{definition}

\begin{definition}[Weak perturbation around unloaded comparison geometry]
In the represented weak-field regime, write
\[
        q^{\rm cmp}_{ab}=\eta_{ab}+h_{ab},
        \qquad |h_{ab}|\ll1,
\]
where $\eta_{ab}$ is the unloaded Lorentzian comparison kernel and $h_{ab}$ the recruited comparison distortion.
\end{definition}

\begin{principle}[Universal comparison-load coupling]
Matter distinctions couple to the represented comparison geometry through the burden they impose on the substrate.  In the tensorial regime this burden is summarized by a stress-energy functional $T^{ab}$, and the leading coupling has schematic form
\[
        \Erec[h;T]
        =\Egeom[h]-\lambda\int h_{ab}T^{ab}\,d^4x
        +O(h^2T,hT^2).
\]
Universality of free fall requires the same comparison kernel to be read by all admissible matter sectors.  The positive cost geometry must be compatible with this coupling through the anchor and cost--causal bridge.
\end{principle}

\begin{gate}[No composition dependence gate]
The equivalence-principle-like claim closes only if the recruitment load of a composite body depends on its total carried stress-energy distinction and not on microscopic composition except through already-accounted internal energy.  All admissible clocks and test bodies must read the same signed comparison kernel and its calibrated cost response.
\end{gate}

\begin{proposition}[Redshift sector as the $h_{00}$ projection]
In the weak static regime, redshift fixes
\[
        h_{00}(r)\simeq\frac{2\Phi_C(r)}{c^2}
        =-\frac{2GM}{rc^2}
\]
up to the sign convention of the comparison kernel.  It does not fix $h_{ij}$, $h_{0i}$, or the positive cost metric $\Gcost$ away from the calibrated clock direction.
\end{proposition}

\begin{proof}
Static redshift compares rest-frame time increments and therefore projects the signed comparison kernel onto its time-time component.  Spatial, mixed, and independent cost-geometry components require additional probes and bridge laws.
\end{proof}

\section{Translation route to the Paper 6 field equations}

The Paper~6 route to GR is not to guess the Schwarzschild metric, but to characterize the admissible local cost functional for \(Q_{ab}\) after the represented smooth Lorentzian gates close.  The expected structure is:
\[
        \Egeom[Q]
        =\int_{\Mgeom} \Bigl[\text{local scalar made from }Q\text{ and its first/second derivatives}\Bigr]\,d\mu_Q.
\]
Paper~6 supplies this represented-geometry closure using H1--H3, A9, and Lovelock uniqueness.  Papers~24/25 already discharge H1 and H2 (continuum emergence and locality of the $\Delta_{\mathrm{geo}}$ closure) at the recruitment level, so the continuum hypotheses needing translation reduce to H3 and the A9/Lovelock step.  The Paper~9 challenge is narrower: translate those residual hypotheses into finite-continuability certificates and identify their weak-field scalar face.

\begin{gate}[Coordinate-relabeling invariance gate]
The cost assigned to a comparison-continuation configuration may not depend on arbitrary labels used for comparison sites.  In the smooth represented regime, this becomes diffeomorphism-type invariance of the geometric cost functional.
\end{gate}

\begin{gate}[Second-order locality gate]
If the substrate cost is local and the represented field equations are to avoid additional higher-derivative propagating comparison modes, the continuum Euler-Lagrange equations should be second order in the metric kernel.  This is the gate that points toward curvature scalar functionals rather than arbitrary higher-curvature actions.
\end{gate}

\begin{gate}[Bianchi/ledger gate]
The geometric side of the represented field equation must have an identically conserved ledger divergence so that local stress-energy continuation is compatible with the geometry equation.  In GR this is the contracted Bianchi identity.  In cost language it is the statement that relabeling comparison sites cannot create or destroy committed load.
\end{gate}


\begin{definition}[Geometric Euler ledger]
Let \(\Egeom[Q]\) be a represented geometric cost functional.  Its Euler ledger is the variational tensor density \(\mathcal E^{ab}\) defined by
\[
        \delta \Egeom[Q]
        =\int \mathcal E^{ab}\,\delta Q_{ab}\,d\mu_Q
        +\text{boundary terms}.
\]
A matter recruitment functional contributes a ledger \(\mathcal T^{ab}\) in the same variational channel.
\end{definition}

\begin{proposition}[Relabeling invariance implies a conserved geometric ledger]
\label{prop:ledger-identity}
If \(\Egeom[Q]\) is invariant under compactly supported infinitesimal relabelings of comparison sites, then its Euler ledger satisfies a formal conservation identity
\[
        \nabla_a\mathcal E^{ab}=0
\]
in the represented smooth metric regime, up to the boundary conventions of the functional.
\end{proposition}

\begin{proof}[Proof sketch]
An infinitesimal relabeling generated by a vector field \(\xi^a\) changes the metric kernel by its Lie derivative,
\[
        \delta Q_{ab}=\Lie_\xi Q_{ab}=\nabla_a\xi_b+\nabla_b\xi_a.
\]
Relabeling invariance gives \(\delta\Egeom=0\) for all compactly supported \(\xi\).  Substituting the variation and integrating by parts yields
\[
        0=-2\int (\nabla_a\mathcal E^{ab})\xi_b\,d\mu_Q.
\]
Since \(\xi\) is arbitrary, \(\nabla_a\mathcal E^{ab}=0\).  The result is a represented-regime conservation identity; translating it back to finite continuations is a separate APF gate.
\end{proof}

\begin{corollary}[Matter compatibility condition]
If the represented field equation has the form
\[
        \mathcal E^{ab}=\lambda \mathcal T^{ab},
\]
then relabeling invariance of the geometric ledger requires
\[
        \nabla_a\mathcal T^{ab}=0
\]
inside the same represented regime.  In APF terms, local committed recruitment load must continue consistently with the comparison geometry that carries it.
\end{corollary}

\begin{theorem}[Translation target: Einstein-Hilbert uniqueness, schematic]
\label{thm:target-eh}
If the represented geometric cost functional is local, smooth, coordinate-relabeling invariant, built only from the tensorial comparison kernel and its derivatives, yields second-order equations, and has a conserved ledger divergence compatible with universal stress-energy recruitment, then the leading continuum functional is of Einstein-Hilbert type,
\[
        \Egeom[Q]\propto \int (R[Q]-2\Lambda)\,d\mu_Q,
\]
up to boundary terms, normalization, and higher-order corrections outside the gate.
\end{theorem}

\begin{warning}[Status of \Cref{thm:target-eh}]
This theorem is not reproved here.  It is the Paper~9 translation target corresponding to the Paper~6 A9/Lovelock closure.  Paper~6 supplies the represented-geometry proof; Paper~9 asks which finite-continuability certificates force the same continuum hypotheses and, in the weak-field scalar sector, which certificate forces the no-trace condition \(\theta=0\).
\end{warning}

\section{Nonlinear alignment: from response law to the Paper 6 closure}

The leading-order theory uses a scalar accumulated load \(\chi\).  Paper~6 supplies the represented tensorial, nonlinear, self-coupled closure to the Einstein equations.  The remaining Paper~9 problem is to show how the finite comparison-substrate response lifts into exactly the Paper~6 tensorial closure rather than an alternative scalar-tensor or higher-derivative geometry.

\begin{gate}[Tensorial geometric response gate]
The scalar stiffness \(K_t\) must lift to a tensorial comparison-cost structure \(g_{ab}\) or equivalent cost kernel such that matter recruitment load changes the full comparison geometry, not only the static time-comparison sector.
\end{gate}

\begin{gate}[Self-coupling gate]
The geometric recruitment load must include the cost of the geometric substrate's own recruited configuration.  In GR language, gravitational field energy gravitates.  In cost language, committed comparison-cost distortion contributes to the future comparison-cost kernel.
\end{gate}

\begin{gate}[Einstein-equation gate]
The nonlinear geometric recruitment functional must have Euler-Lagrange equations equivalent, in the represented smooth Lorentzian regime, to
\[
        G_{ab}+\Lambda g_{ab}=\frac{8\pi G}{c^4}T_{ab},
\]
with \(G\) remaining either a substrate-specific input or a derived stiffness constant.
\end{gate}

\begin{warning}[Primary translation theorem]
The Paper~6 endpoint can be restated in Paper~9 language as follows:

\emph{If finite comparison continuations admit a complete calibrated response geometry, a physical comparison-response anchor, a smooth local Lorentzian signed comparison kernel, universal recruitment coupling, local A9-type closure, and nonlinear self-consistent response, then their represented continuum equations are the Einstein equations.}

Paper~6 proves the represented-geometry closure under its hypotheses.  Paper~9 has not yet proved that finite continuability forces those hypotheses.  Its current hard local target is the weak-field trace theorem \(\theta=0\), the scalar shadow of the same closure.
\end{warning}


\begin{audit}[Current cross-corpus status]
\begin{longtable}{p{0.31\textwidth}p{0.12\textwidth}p{0.46\textwidth}}
\toprule
Object & Status & Role here \\
\midrule
Held continuation profile & I & Imported from Papers~1, 10, and 14 as the pre-record identity carrier. \\
Record-forming class transition & I & Imported from Paper~37; downstream of the held response geometry. \\
Complete represented profile carrier & I & Imported from Papers~1 and 10; Paper~5 supplies the quantum-sector completion. \\
Local quantum orientation & I, quantum only & Imported from Paper~5 only in its declared coherent quantum sector; not identified with clock--ruler or causal orientation. \\
Geometric response quotient & P$\mid$H & Proved here under constant-rank regularity. \\
Substrate anchor & O/C & Identity-complete physical anchor still to be constructed. \\
Lorentzian signed comparison kernel & O/C & Requires independent cone and signed two-homogeneous gates. \\
No-trace weak-field source & O/C & The invariant target remains $\Ptr\mathcal K^{-1}\mathcal J=0$. \\
\bottomrule
\end{longtable}
\end{audit}

\section{Input ledger and closure plan}

\begin{longtable}{@{}p{0.265\textwidth}p{0.115\textwidth}p{0.515\textwidth}@{}}
\toprule
Claim / object & Status & What remains \\
\midrule
Continuation judgment $\G\entails_C D\goes E$ & \Status{P} & Imported from Paper~10 as primitive grammar. \\
Finite operational basis & \Status{R/G+C} & General theorem requires joint extension, minimum attainment, closure invariance, and positive marginal extension cost; finite atom-cover model is Engine-certified. \\
Operational basis to physical precision channels & \Status{O/G} & Construct the faithful realization preserving joint realizability and recoverability, separating reduced states, and admitting real limits. \\
Record topology and second countability & \Status{P$\mid$G} & Hausdorffness requires state-separating real channels; second countability requires the countable grammar and topology-generation gate. \\
Finite smooth generation & \Status{O/G} & Prove local smooth functional generation by the finite channel map. \\
Smooth response-factor completeness & \Status{O/G} & Prove smooth two-way factorization between the complete response and a finite local APF generator before using the derivative-kernel identity. \\
Regular response quotient & \Status{P$\mid$H} & Proved locally after constant rank and smooth response-factor completeness; the physical comparison response model and rank stratification remain to be constructed. \\
Response-orbit gauge and distance & \Status{P} & Proved intrinsically on the regular quotient; smooth strongly convex Finsler structure requires additional norm regularity. \\
Linear Paper~2 specialization & \Status{P} & Proved as the flat quotient branch with optimal response gains. \\
Physical-ledger calibration & \Status{C/P} & Comparison theorem proved; actual ledger-rate bounds and response-class completeness remain physical inputs. \\
Residual energy--distance relation & \Status{P$\mid$SRR} & Proved locally when the residual Hessian is positive definite. \\
Physical comparison-substrate anchor & \Status{O/C} & Construct $s$ and $A=q\circ Ds$, type the null kernel and provenance, and prove direct-germ calibration. \\
Quadratic positive cost metric & \Status{P$\mid$Q} & Proved from the parallelogram certificate; whether the substrate passes it is open. \\
Signed Lorentz--Finsler comparison structure & \Status{O/C} & Derive the causal cone, orientation, signed two-homogeneous function, and Lorentzian Hessian. \\
Dimension lock-in & \Status{O} & Requires an anchor-aware mode-count theorem, not merely register count. \\
Point-source $1/r$ profile & \Status{C/P} & Proved from an assumed isotropic quadratic exterior kernel. \\
Clock response and redshift & \Status{R} & Conditional recovery under load law, calibration, and proper-time dictionary. \\
Trace--volume identity & \Status{P/C} & Proved in the represented clock--ruler comparison plane. \\
Gauge-fixed response Hessian & \Status{O/C} & Derive $\mathcal K$ and $\mathcal J$ from the recruitment master equation after gauge removal. \\
No-trace operator $\Ptr\mathcal K^{-1}\mathcal J=0$ & \Status{O/C} & Decisive local theorem; source-selection and fixed-volume are the principal routes. \\
Spatial response $\gamma_C=1$ & \Status{C/O} & Follows from the no-trace operator through the trace--volume theorem. \\
Conditional weak-field exterior & \Status{C/R} & Follows from clock calibration, universal coupling, and no trace. \\
Full Einstein equations & \Status{I/C} & External Paper~6 endpoint; residual finite-continuability translation remains open. \\
\bottomrule
\end{longtable}

\section{Formal theorem tree for closure}

\begin{longtable}{p{0.10\textwidth}p{0.30\textwidth}p{0.11\textwidth}p{0.34\textwidth}}
\toprule
ID & Claim & Status & Dependency \\
\midrule
G9.1 & A finite complete operational basis exists. & \Status{R$\mid$G/C} & General joint-extension package; exact atom-cover Engine witness. \\
G9.2 & The operational basis is realized as finite separating precision channels. & \Status{O/G} & Faithful operational-to-channel realization. \\
G9.3 & The record topology is Hausdorff and second countable. & \Status{P$\mid$G} & State separation; countable grammar; topology-generation gate. \\
G9.4 & Finite smooth generation gives a subcartesian differential space. & \Status{P$\mid$G} & Finite smooth-generation gate. \\
G9.5 & Complete local continuation flows generate smooth orbits. & \Status{I/P} & Subcartesian orbit theorem. \\
G9.6 & Smooth response-factor completeness gives $\ker D\beta=N^{\rm ctx}$. & \Status{P$\mid$G} & Two-way smooth factorization gate. \\
G9.7 & The process norm defines an intrinsic response-orbit quotient gauge and length. & \Status{P} & \Cref{thm:canonical-operational-anchor,thm:response-orbit-length}. \\
G9.8 & A physical anchor pulls response-orbit length back to site geometry. & \Status{P$\mid$G} & \Cref{gate:physical-substrate-anchor,thm:anchored-cost-geometry}. \\
G9.9 & Calibrated physical path cost is bi-Lipschitz to response distance. & \Status{P$\mid$H} & \Cref{thm:ledger-calibration}. \\
G9.10 & Smooth residual energy is one half squared local distance. & \Status{P$\mid$SRR} & \Cref{thm:residual-energy-distance}. \\
G9.11 & Parallelogram quadraticity gives a positive cost metric. & \Status{P$\mid$Q} & \Cref{prop:positive-cost-metric}. \\
G9.12 & Positive pullback geometry cannot supply Lorentzian signature. & \Status{P} & \Cref{prop:positive-pullback-no-lorentz}. \\
G9.13 & A signed Lorentz--Finsler comparison structure represents the causal cone. & \Status{O/C} & \Cref{gate:signed-comparison}. \\
G9.14 & Isotropic quadratic exterior recruitment gives $\chi\sim1/r$. & \Status{C/P} & \Cref{prop:point-source-profile}. \\
G9.15 & Calibrated clock response conditionally recovers redshift. & \Status{R} & Response law and proper-time dictionary. \\
G9.16 & Generator trace equals comparison-volume variation. & \Status{P/C} & \Cref{thm:trace-volume}. \\
G9.17 & Linear loaded equilibrium is $h_\chi=\chi\mathcal K^{-1}\mathcal J+O(\chi^2)$. & \Status{P$\mid$H} & \Cref{thm:no-trace-operator}. \\
G9.18 & The no-trace gate is $\Ptr\mathcal K^{-1}\mathcal J=0$. & \Status{O/C} & Source-selection or constrained-volume theorem. \\
G9.19 & No trace gives $\gamma_C=1$ and the first-PN exterior. & \Status{C/R} & \Cref{cor:ppn-volume-mode,thm:conditional-weak-field}. \\
G9.20 & The represented nonlinear closure is Einstein gravity. & \Status{I/C} & External Paper~6 H1--H3/A9/Lovelock endpoint. \\
\bottomrule
\end{longtable}

\section{Minimal publishable core}
\label{sec:minimal-publishable-core}

A stable geometry-centered version of Paper~9 should contain:
\begin{enumerate}
\item the regular response quotient and pullback-gauge theorem;
\item the intrinsic path-distance and linear Paper~2 recovery theorem;
\item the physical-ledger calibration theorem and the length--energy fork;
\item the intrinsic response-orbit gauge, the physical comparison-anchor gate and conditional pullback theorem, and the explicit separation of positive cost from Lorentzian comparison geometry;
\item the trace--volume theorem in the loaded clock--ruler plane;
\item the no-trace operator reduction $\Ptr\mathcal K^{-1}\mathcal J=0$;
\item the existing conditional redshift and first-PN recovery package, with Paper~6 imports kept explicit.
\end{enumerate}
This is a stronger mathematical identity for the paper than ``cost accounting plus a weak-field bridge.''  It makes Paper~9 the intrinsic differential geometry of admissibility response while preserving the precise gravitational target.

\begin{audit}[Defer connections and curvature]
A connection is not canonical on a general nonsmooth or non-strongly-convex Finsler structure, and rank-changing response spaces may not be manifolds.  Curvature should therefore be introduced only after the regular quotient, smoothness, strong convexity or quadraticity, anchor, and causal signature gates have closed.  The trace--volume theorem does not require that later machinery.
\end{audit}

\section{Research tasks for the next draft}

\begin{enumerate}
\item \textbf{Construct the physical response manifold.}  Identify the actual comparison-response coordinates in the recruitment formalism, compute the rank of the response map, and classify rank-changing strata.
\item \textbf{Calibrate the ledger rate.}  Prove bounds $a\Fresp\le j\le b\Fresp$ on a complete declared response class, or exhibit fixed setup, threshold, or hysteresis terms that force a non-length branch.
\item \textbf{Derive the comparison-response anchor.}  Specify how an infinitesimal site displacement recruits a quotient response direction and prove agreement with the direct comparison germ.
\item \textbf{Test quadraticity rather than assume it.}  Measure or derive parallelogram defects.  Retain the Finsler branch if weighted or polyhedral cost survives.
\item \textbf{Build the causal comparison function.}  Derive the finite-coordination cone, irreversible orientation, and a signed two-homogeneous $\Qcmp$ with Lorentzian Hessian.
\item \textbf{Compute the loaded response operator.}  Expand the Papers~24/25 recruitment master equation in the comparison-kernel sector to identify $\mathcal K$ and $\mathcal J$ after gauge removal.
\item \textbf{Prove the no-trace operator.}  Pursue either $\Ptr\mathcal J=0$ plus Hessian block-diagonality, or a primitive fixed-volume constraint.  The target is $\Ptr\mathcal K^{-1}\mathcal J=0$, not merely a fitted $\gamma_C$.
\item \textbf{Interpret the weight-one program microscopically.}  Treat its capacity accounting as a candidate mechanism for source-selection, not as a substitute for the general response theorem.
\item \textbf{Freeze Paper~2.}  Once these theorems stabilize, Paper~2 should import the finite linear specialization and its physical-ledger bridge rather than develop a parallel geometry section.
\end{enumerate}

\appendix

\section{Candidate microscopic realization of the no-trace gate: the weight-one program}
\label{sec:weight-one-reduction}

The trace gate can be priced in capacity terms, and the price splits.  Write the static response as a direct clock coupling plus a spread deposit of weight \(w\), distributed \(\eta\)-democratically over the comparison cell with denominator \(n-1\), where \(n\) is the spatial dimension.  We adopt this decomposition and this spread shape; neither is derived, and the ledger below carries them as S0 and S1.  The spatial coefficient is then a one-parameter family,
\[
        \gamma_C(w)=\frac{w}{(n-1)-w},
\]
and general relativity is the weight-one point in every dimension: \(\gamma_C(1)=1/(n-2)=\gamma_{\rm GR}(n)\).  The closure candidates for the response sit at fork coordinates on this curve --- the plane rule at \(w=1\), the four-cell rule at \(w=1/2\), the per-channel rule at \(w=0\), the pooled rule at no finite weight --- so the gate question becomes a question about one scalar weight \coderef{check\_T\_gammaC\_carrier\_fork}{gamma\_c\_carrier\_program.py}.

The weight has an exact accounting reading.  Let \(D\) be the cell-sum of a deposit pattern over the \(n+1\) channel accounts of the comparison cell.  Then \(w=D(\text{spread deposit})/D(\text{commitment})\), identically in \(n\) and \(w\): the \(\eta\)-diagonal sums to \(n-1\) over the cell, so the denominator is not an independent import but the unique normalization under which the weight reads as a fraction of the committed amount.  The same cell-sum annihilates every clock--ruler exchange pattern, so the redshift sector carries zero net amount and cannot be over-billed by this accounting.  And \(w\neq1\) is exactly one identically conserved scalar improvement mode of strength proportional to \(1-w\) --- the independent scalar comparison mode of \Cref{prop:nonzero-trace-scalar-mode}, met in ledger coordinates \coderef{check\_T\_weight\_one\_reduction}{gamma\_c\_carrier\_program.py}.  In this coordinate the no-trace gate is a conservation statement: the substrate's cell-side deposit matches the load's committed amount.

Matching splits into two halves with different prices.  The no-source half needs the minimal trace-currency premise TCP-min --- the response deposit is denominated in committed capacity at a channel-uniform rate --- a named conjecture, not a theorem.  Under it the deposit must be billed to the load's own committed line: no census sector furnishes a second local payer, and a line cannot be overdrawn.  Hence \(D(\text{deposit})\le D(\text{commitment})\), that is \(w\le1\) and \(\gamma_C\le1/(n-2)\); at \(n=3\),
\[
\begin{split}
        \gamma_C\le1
        \quad\text{at}\quad
        [\,&P_{\rm structural}\mid \text{S0/S1}+\text{TCP-min}+\text{H5-W}+\text{census legs}\\
        &+\text{the Paper 9/Paper 10 base (P5, B7 incl.\ adopted \(C_M\))}\,].
\end{split}
\]
The bound is citable only with that ledger attached.  H5-W is census transport re-fenced to net-deposit objects --- a named premise, separate from the rest and still standing on its own.  P5 is this paper's scalar package; B7 is Paper~10's held-record capacity floor together with the adopted identification of the maintenance cost \(C_M\) with committed capacity.  The no-sink half costs more.  It needs TCP-strong, the re-expression-in-full clause: the static response is the same committed amount read at the cell, not less --- the one-sided clause \(D(\text{deposit})\ge D(\text{commitment})\).  Nothing cheaper closes it, because an absence of deposit separates no continuation profile and carries no bill.  Composed with the bound, TCP-strong gives \(w=1\), hence \(\gamma_C=1/(n-2)\) in every \(n\) and \(\gamma_C=1\) at \(n=3\) --- still conditional on the full ledger.  Neither trace-currency clause is derived.

The TCP-strong gap has a precise location in Paper~10.  Testedness there is binary: a distinction-expression either separates a tested continuation profile or preserves every one.  The amount-graded refinement that would close the gap --- billing a partial expression for its missing fraction --- does not derive from the calculus.  The obstruction is an encoding wall: the calculus grades the tested-atom count, the count is representation data, and floor-coherent signatures hold the same record at grades \(\eps^\ast\cdot\{1,\log_2N,N\}\) depending on how it is atomized, with \(\eps^\ast\) the distinction floor \coderef{check\_L\_amount\_graded\_testedness\_encoding\_no\_go}{amount\_graded\_testedness.py}.  What survives is an anatomy: over the adopted linear frame the discharge splits into a distributed-encoding existence premise and a par billing rate, with the dimension audit selecting par within the linear family --- a selection, not a discharge.  Whether the substrate holds a static response profile as one compressed register or as one tested unit per resolution rung is a record-structure question about the substrate, and it is the single genuine reopening door.

The door has a shape.  In the coherent shared-read model, the only class of activation law that moves the encoding argmin toward the distributed register is a strictly superadditive read-contention law; linear and subadditive laws leave the compressed register in place \coderef{check\_T\_fagt2\_encoding\_argmin\_pressure\_and\_read\_channel}{amount\_graded\_testedness.py}.  And the banked corpus prices every compatible-sharing configuration on the shut side of that gate.  The derived joint-cost identity \(\Delta=\eps^\ast(J-R)\), iterated over \(M\) compatible readers of one shared record channel, prices the sharing at \(\Delta=-(M-1)\,\eps^\ast\) --- a strictly deepening discount; the banked activation floor is linear and saves nothing; and the recruitment corpus is reader-blind by census, with no demand-multiplicity variable anywhere in its banked signatures.  The one sign-positive source in the calculus is irreducible jointness, \(J>R\) --- the occupancy feature.  So the residual question is an occupancy fork.  Either the concurrent-activation bottleneck sits at sub-record granularity --- finiteness supplies the superadditive shape of such a cap for free, but the granularity is a premise finiteness does not force --- or concurrent readouts of the static record are irreducibly joint at atom granularity: occupancy at the read interface, in the current Paper~5 sense of a non-Boolean held-profile sector with record-incompatible co-available continuations, per-interface and read off the world rather than derived, and not adopted here.  The bill for the second horn is stated plainly: it needs the static-field read interface to present record-incompatible co-available families, and a classical broadcast readout of one static record is the separable configuration par excellence, so the natural typing sits on the adverse side \coderef{check\_T\_contention\_law\_granularity\_occupancy\_fork}{amount\_graded\_testedness.py}.

The bound meets observation signed.  Cassini (Bertotti--Iess--Tortora 2003) gives \(\gamma-1=(2.1\pm2.3)\times10^{-5}\), and the central value sits \(0.91\sigma\) \emph{above} the bound: the bound survives at under one sigma and is under live one-sided tension, not merely saturated.  A confirmed \(\gamma_{\rm obs}>1\) at \(\ge3\sigma\) falsifies the composite (S0/S1 + TCP-min + H5-W + census + P5/B7 base).  A factor-three precision improvement at the same central value puts the composite under roughly \(3\sigma\) pressure.  A confirmed \(\gamma_{\rm obs}<1\) at \(\ge3\sigma\) refutes only the no-sink clause and leaves the bound standing.  The two halves are separately falsifiable, which the unsplit gate was not.

The decisive target remains the no-trace/unimodular certificate of \Cref{gate:unimodular-load-response}: by \Cref{prop:unimodular-implies-gamma,cor:spatial-target-equivalence} it closes the leading spatial coefficient, and by \Cref{prop:generator-trace-obstruction,prop:trace-obstruction-gamma} any trace source is exactly the PPN-\(\gamma\) obstruction.  What the reduction contributes is the exact residue: over the adopted spread shape, the gate's underived content is two currency clauses, and the cheaper one already buys a falsifiable one-sided bound.


\section{Source-of-record redshift calculation}

This appendix records the compact calculation that motivates the first recovery theorem.  Let a static point mass carry an effective accumulated recruitment load \(\chi_{\rm rec}(r)\).  Suppose the exterior linear-response law is
\[
        K_t(r)=K_{t0}(1-\alpha\chi_{\rm rec}(r)),
\]
and suppose the point-source equilibrium plus physical calibration gives
\[
        \alpha\chi_{\rm rec}(r)=\frac{2GM}{rc^2}.
\]
Then
\[
        K_t(r)=K_{t0}\left(1-\frac{2GM}{rc^2}\right).
\]
Using the represented proper-time dictionary,
\[
        d\tau(r)=\sqrt{K_t(r)/K_{t0}}\,dt,
\]
one obtains
\[
        d\tau(r)=\sqrt{1-\frac{2GM}{rc^2}}\,dt.
\]
For two static observers at \(r_B<r_A\), equal coordinate periods correspond to different local proper periods, hence
\[
        \frac{\omega_A}{\omega_B}
        =\frac{d\tau(r_B)}{d\tau(r_A)}
        =\sqrt{\frac{1-2GM/r_Bc^2}{1-2GM/r_Ac^2}}.
\]
Expanding to first order gives
\[
        \frac{\omega_A}{\omega_B}
        \approx
        1-\frac{GM}{c^2}\left(\frac1{r_B}-\frac1{r_A}\right).
\]
This is the static gravitational redshift form.  The calculation is a conditional recovery because the response law, calibration, and proper-time dictionary remain open gates.

\section{Minimal next-lemma inventory}

The next version should prove, weaken, or abandon the following lemmas.

\begin{enumerate}
\item \textbf{Regular-quotient lemma.}  The comparison-response map supplied by the recruitment formalism has constant rank on the weak-field branch, or else admits a controlled stratification.
\item \textbf{Ledger-calibration lemma.}  The physical interaction rate is uniformly comparable with the pullback process gauge on the complete response class.
\item \textbf{Anchor lemma.}  The direct comparison germ equals the anchored response germ for all locally admissible comparison directions.
\item \textbf{Quadraticity-or-Finsler lemma.}  Determine whether the local gauge passes the parallelogram certificate; if not, retain the exact Finsler geometry rather than fitting a metric.
\item \textbf{Causal-cone lemma.}  Finite coordination and irreversible record order produce a smooth oriented cone and signed comparison function with Lorentzian Hessian.
\item \textbf{Response-Hessian lemma.}  Loaded-equilibrium expansion yields a gauge-fixed invertible $\mathcal K$ and source $\mathcal J$ in the comparison-kernel sector.
\item \textbf{No-trace source lemma.}  Prove $\Ptr\mathcal J=0$ and $\Ptr\mathcal K\Ptf=0$, or derive an admissibility constraint enforcing $\operatorname{tr}_{q_0}h=0$.
\item \textbf{Self-coupling lemma.}  The cost of geometric distortion contributes to later recruitment load and lifts the weak-field response into the Paper~6 nonlinear closure.
\item \textbf{Ledger-divergence lemma.}  Relabeling comparison sites induces an identically conserved geometric ledger compatible with matter continuation.
\end{enumerate}

\section{Conclusion}

Paper~9 is now organized around a precise geometric claim.  Continuation is primitive, while geometry is represented continuation.  A structurally aligned held continuation profile is first represented in a complete carrier and process algebra.  A finite response model then loses its redundant implementation coordinates by quotienting response fibers.  The represented process norm then pulls back to a positive tangent gauge, whose path integral gives intrinsic response distance.  In a linear chart this geometry is flat and reproduces the quotient norm, optimal upper gain, positive minimum gain, and compatible bi-Lipschitz comparison already established in Paper~2.  A separately calibrated physical ledger may equal that distance, remain only comparable with it, or fail to be a path metric because of setup, threshold, indivisibility, saturation, or hysteresis.

The same architecture resolves the homogeneity ambiguity.  One-homogeneous path ledgers generate length.  Smooth residual penalties generate local quadratic energy, and their correct geometric form is one half squared distance.  Neither branch should be renamed into the other.  A conditional operational basis, followed by a faithful physical-channel realization, countable topology generation, and finite smooth generation, yields the record differential-space branch on which complete continuation flows form immersed smooth orbits.  Smooth response-factor completeness then licenses the derivative-kernel identity and the intrinsic response-orbit quotient gauge.  A physical comparison-substrate anchor is not obtained merely by naming that quotient: it remains a separate site-map and tangent-compatibility gate.  Under quadraticity the intrinsic gauge yields a positive cost metric, while a separate causal cone and signed two-homogeneous structure are required for Lorentzian comparison geometry.  This blocks the shortcut $g=T^*T$ as a derivation of spacetime signature while preserving the valid conditional pullback construction for positive cost.

The gravitational application becomes sharper rather than weaker.  Recruitment load deforms the represented clock--ruler comparison kernel.  The generator trace is exactly the first variation of the local comparison-volume element, and clock calibration turns that scalar volume mode into the PPN obstruction $1-\gamma_C$.  In the gauge-fixed loaded-equilibrium problem the entire no-trace question is
\[
        \Ptr\mathcal K^{-1}\mathcal J=0.
\]
This equation separates two routes: show the scalar source is trace-free and the Hessian does not mix trace sectors, or derive a primitive fixed-volume admissibility constraint.  The weight-one capacity program may supply a microscopic source-selection mechanism, but the invariant operator equation is the theorem it must discharge.

Paper~6 remains the represented smooth source of record for the Einstein closure.  Paper~9 supplies the intrinsic response geometry underneath it: quotient, response-orbit gauge, path distance, ledger-calibration theorem, physical-anchor gate and conditional pullback, causal-signature gate, recruitment deformation, and trace--volume response.  That is the coherent role of the paper and the correct home of the interaction-geometry program.

\begin{quote}
Paper~9 establishes a dependency-explicit route from finite operational comparison content to record topology, Sikorski differential structure, subcartesian regular loci, smooth continuation orbits, and positive response-orbit geometry. The general finite-basis theorem is conditional on a named joint-extension package and is machine-witnessed exactly on the finite atom-cover model. Physical precision-channel generation, finite smooth generation, smooth response-factor completeness, the physical substrate anchor, Lorentzian signature, the physical no-trace source law, and the nonlinear Einstein endpoint remain separately typed gates. The weak-field trace reduction is therefore a conditional bridge, not an unconditional rederivation of general relativity.
\end{quote}

\begin{center}
\fbox{\begin{minipage}{0.90\textwidth}
\centering
\vspace{3pt}
\textbf{Compressed result.}\\[3pt]
Held continuation profiles supply the pre-record identity carrier.  Response-null quotienting of their represented process realizations produces intrinsic process-induced geometry.  Physical cost equals that geometry only after ledger calibration.  The comparison substrate inherits it only through an anchor.  Quadraticity yields a positive cost metric; causal signature requires a separate signed gate.  Under recruitment load, the scalar weak-field obstruction is the local comparison-volume mode
\[
        \operatorname{tr}L
        =\left.\frac{d}{d\chi}\log\sqrt{|\det q_\chi|}\right|_0
        =1-\gamma_C,
\]
and its decisive microscopic target is $\Ptr\mathcal K^{-1}\mathcal J=0$.\\[3pt]
\end{minipage}}
\end{center}

\end{document}
